%%%%%%%% ICML 2026 SUBMISSION FILE %%%%%%%%%%%%%%%%%

\documentclass{article}

% Recommended packages
\usepackage{microtype}
\usepackage{graphicx}
\usepackage{booktabs}
\usepackage{hyperref}
\usepackage{xcolor}

\newcommand{\theHalgorithm}{\arabic{algorithm}}

% For blind submission:
\usepackage{icml2025}  % Update to icml2026 when available
% For camera-ready: \usepackage[accepted]{icml2025}

\usepackage{amsmath}
\usepackage{amssymb}
\usepackage{amsthm}
\usepackage{algorithm}
\usepackage{algorithmic}
\usepackage{enumitem}

% Theorems
\theoremstyle{plain}
\newtheorem{theorem}{Theorem}[section]
\newtheorem{proposition}[theorem]{Proposition}
\newtheorem{lemma}[theorem]{Lemma}
\newtheorem{corollary}[theorem]{Corollary}
\theoremstyle{definition}
\newtheorem{definition}[theorem]{Definition}
\newtheorem{remark}{Remark}

\icmltitlerunning{STED: Semantic Similarity for Complex JSON Structures}

\begin{document}

\twocolumn[
\icmltitle{STED: Semantic Tree Edit Distance for JSON Similarity \\with Application to LLM Consistency Evaluation}

\icmlsetsymbol{equal}{*}

\begin{icmlauthorlist}
\icmlauthor{Anonymous Author(s)}{anon}
\end{icmlauthorlist}

\icmlaffiliation{anon}{Anonymous Institution}

\icmlcorrespondingauthor{Anonymous Author(s)}{anonymous@example.com}

\vskip 0.3in
]

\printAffiliationsAndNotice{\icmlEqualContribution}

\begin{abstract}
Large Language Models (LLMs) are increasingly deployed for structured data generation, yet output consistency remains critical for production applications. We present \textbf{STED (Semantic Tree Edit Distance)}, a similarity metric for complex JSON structures that combines structural tree comparison with semantic understanding via embeddings and order-invariant matching. Applying STED to evaluate LLM consistency through \textbf{large-scale benchmarking} of 18 LLMs across 10 providers (2.2M+ outputs), we reveal: (1) dramatic task-specific gaps---Llama-3.3-70B achieves 99\% structured output validity but only 33\% on tool calling; (2) temperature strongly affects stability---consistency drops 22--57\% from T=0 to T=1; (3) complex parameter types strongly predict inconsistency. Our \textbf{consistency diagnostics framework} identifies parameter complexity as the primary predictor (27--29$\times$ higher for lists/dicts across all temperatures; 60$\times$ for deeply nested structures at T=0.5), with four distinct failure patterns enabling targeted interventions. We introduce stability score ($S_\alpha$) for efficient model ranking; embedding ablation across 9 models confirms robustness ($r>0.999$). While validated in the LLM context, STED's design principles may generalize to other JSON comparison tasks such as API testing and configuration management.
\end{abstract}

\let\thefootnote\relax\footnotetext{A preliminary version of this work appeared at a NeurIPS 2025 workshop. This submission substantially extends it with: (1) expanded LLM benchmarking (18 models across 10 providers, 2.2M+ outputs), (2) consistency diagnostics framework identifying failure patterns and predictors, (3) comprehensive ablation studies on stability score and embedding models, and (4) practitioner guidelines for deployment.}

\section{Introduction}

Large Language Models (LLMs) are increasingly deployed for structured data generation in JSON format for APIs, data extraction, and automated workflows. However, evaluating consistency in LLM-generated structured outputs presents unique challenges that existing methods fail to address adequately.

The fundamental issue is a mismatch between evaluation methods and structured semantics. BERTScore~\cite{Zhang2020}, while effective for natural language, suffers from order sensitivity that violates JSON's order-agnostic semantics, underestimating similarity for reordered but identical structures. Tree Edit Distance (TED) focuses purely on structure without semantic understanding. DeepDiff~\cite{seperman} can ignore ordering but treats \texttt{"email"} and \texttt{"email\_address"} as completely different, missing obvious semantic relationships.

LLMs frequently generate functionally equivalent structures with semantic variations: naming conventions (\texttt{"user\_name"} vs \texttt{"userName"}), array reorderings, and type representations (\texttt{"true"} vs \texttt{true}). Existing metrics fail to recognize these equivalences, causing false negatives in production systems.

We propose \textbf{STED (Semantic Tree Edit Distance)}, a similarity metric for complex JSON structures that combines: (1) \textbf{Semantic-Enhanced Comparison} recognizing equivalent keys and values via embeddings; (2) \textbf{Order-Invariant Matching} via Hungarian algorithm; (3) \textbf{Multi-Level Similarity} integrating structural, key, value, and type dimensions. While STED integrates established techniques (tree edit distance, optimal assignment, embedding similarity), \textbf{no existing tool combines them for structured data comparison}---this gap leaves practitioners without reliable metrics for JSON similarity measurement. We apply STED to evaluate LLM output consistency; while the metric's design principles are general, our validation focuses on the LLM domain. Our experiments show STED correctly identifies schema violations ($0.23$ similarity) while scoring $>0.95$ for benign reorderings---$4\times$ better discrimination than baselines.

\textbf{Why is this empirical work necessary?} As LLMs increasingly power production APIs and automated workflows, unreliable consistency evaluation leads to: (1) false negatives rejecting valid outputs, (2) false positives accepting breaking changes, and (3) inability to compare models or tune parameters with confidence. We provide the first systematic study addressing this gap:

\begin{itemize}[leftmargin=*, nosep]
\item \textbf{Large-Scale Benchmarking}: We evaluate 18 LLMs across 10 providers on two complementary tasks---structured output generation (75 samples) and tool calling (1,006 samples)---comprising 2.2M+ total outputs with rigorous statistical validation (Friedman tests, bootstrap CIs), revealing critical task-specific variations.

\item \textbf{Consistency Diagnostics Framework}: We provide the first systematic analysis of inconsistency sources, finding parameter complexity strongly predicts inconsistency (27--29$\times$ higher for complex types). We characterize four distinct failure patterns---string drift (39\%), tool count variation (35\%), tool selection (16\%), and complex parameters (6\%)---enabling targeted intervention strategies.

\item \textbf{Human Validation}: A focused study on discriminative cases (74 samples, 3 annotators, Fleiss' $\kappa=0.759$) confirms STED's alignment with human judgment (87.8\% via majority vote), outperforming DeepDiff and TED significantly ($p<0.001$) with a notable gap over BERTScore (13.5pp, $p=0.087$), complementing synthetic validation (Section~\ref{sec:human-validation}).

\item \textbf{Theoretical Characterization}: We formally characterize STED's properties and introduce stability score ($S_\alpha$) for efficient model ranking. Embedding ablation across 9 models confirms robustness ($r>0.999$).
\end{itemize}

\section{Related Work}

\textbf{Tree Edit Distance.} The tree edit distance problem, introduced by Tai~\cite{Tai1979} and made practical by Zhang-Shasha~\cite{Zhang1989}, provides theoretical foundations for comparing hierarchical structures through insert, delete, and update operations. Subsequent work improved computational efficiency: RTED~\cite{Pawlik2011} achieves optimal worst-case complexity, while learned approaches~\cite{paassen2018tree,paassen2022revisiting} adapt costs to specific domains. For JSON comparison, tools like DeepDiff~\cite{seperman} offer order-insensitive matching via iterative deepening, but lack semantic awareness---treating \texttt{"email"} and \texttt{"email\_address"} as completely unrelated despite obvious semantic equivalence. JSONPath~\cite{Gessner2007} and JSON Patch~\cite{RFC6902} focus on querying and modification rather than similarity measurement.

\textbf{Neural Text Similarity.} BERTScore~\cite{Zhang2020} revolutionized text similarity by using contextualized BERT~\cite{Devlin2019} embeddings with greedy token alignment, achieving high correlation with human judgments. However, its position-sensitive matching violates JSON's order-agnostic semantics---reordering array elements or object keys produces different scores despite identical meaning. Sentence-level approaches like Sentence-BERT~\cite{Reimers2019}, SimCSE~\cite{Gao2021}, and Universal Sentence Encoder~\cite{Cer2018} process JSON as flat strings, losing hierarchical structure. Graph neural networks~\cite{Hamilton2017,Velickovic2018} and Tree-LSTMs~\cite{Tai2015} can model structure but require task-specific training and are computationally expensive for online comparison.

\textbf{LLM Structured Output and Consistency.} The emergence of LLMs has sparked interest in structured generation~\cite{Bubeck2023,Liang2022} and evaluation frameworks~\cite{yang2025structeval,shorten2024structuredrag}. Recent work on format control includes grammar-constrained decoding~\cite{wang-etal-2025-verifiable} and JSON schema enforcement~\cite{geng2025jsonschema}. LLM consistency research has examined paraphrase consistency~\cite{elazar2021measuring}, semantic stability under input perturbations~\cite{Raj2023,atil2024llm}, temperature effects on output variability~\cite{renze2024effect,renze-2024-effect}, and non-determinism in ``deterministic'' settings~\cite{baldwin2024nondeterminism}. However, this body of work focuses primarily on natural language outputs or input-side variations, leaving structured output consistency for identical inputs largely unexplored.

\textbf{Why Not LLM-as-Judge?} While LLM-based evaluation has shown promise for natural language tasks~\cite{zheng2023judging}, it is problematic for consistency measurement: (1) \textit{circular reasoning}---using potentially inconsistent LLMs to judge LLM consistency; (2) \textit{judge variability}---LLMs produce different outputs across runs even at temperature 0, so judge scores would themselves vary, defeating the purpose of consistency measurement; (3) \textit{computational cost}---evaluating 2.2M+ outputs requires millions of pairwise comparisons, infeasible with API calls; (4) \textit{interpretability}---STED produces deterministic scores decomposable into structural and semantic components, unlike black-box LLM judgments. We validate STED against human judgment directly (Section~\ref{sec:human-validation}) rather than through an LLM proxy.

\textbf{Semantic and Schema Matching.} Ontology alignment~\cite{Euzenat2013} and schema matching~\cite{Rahm2001,Do2002,Melnik2002} combine lexical, structural, and constraint-based strategies to identify correspondences between heterogeneous data sources. The Hungarian algorithm~\cite{Kuhn1955,Munkres1957} provides optimal bipartite matching in polynomial time, enabling principled alignment of unordered collections. Entity matching~\cite{Mudgal2018} and change detection~\cite{Chawathe1996} address related problems in databases. We adapt these techniques for LLM output comparison, combining embedding-based semantic similarity with optimal assignment for order-invariant JSON matching.

\textbf{Summary.} Existing methods address either structure (TED, DeepDiff) or semantics (BERTScore, embeddings) but not both; are either order-sensitive (BERTScore) or semantically blind (DeepDiff); and miss common variation patterns in structured data like naming convention differences and type coercion. STED bridges these gaps through unified structural-semantic analysis with order-invariant matching, providing a general-purpose JSON similarity metric that we apply to LLM consistency evaluation.

\section{Methodology}

We present STED (Semantic Tree Edit Distance), a similarity metric for complex structured data that integrates established techniques to address gaps in existing methods. While traditional metrics treat structured formats like JSON, XML, and HTML as flat text, we recognize their hierarchical nature by transforming them into tree representations. STED extends classical tree edit distance with semantic similarity capabilities to compute meaningful similarity scores between any two structured documents. For LLM consistency evaluation, we aggregate pairwise STED scores across multiple generations into normalized consistency metrics. However, STED itself is a general-purpose similarity measure applicable to any structured data comparison task.

\subsection{Problem Formulation}

\textbf{Core Problem: JSON Similarity.} Given two JSON documents represented as trees $T_1$ and $T_2$, compute a similarity score $s(T_1, T_2) \in [0, 1]$ that captures both semantic and structural equivalence, where $s = 1$ indicates identical structures and $s = 0$ indicates fundamental incompatibility. While our framework applies to any hierarchical format (JSON, XML, HTML), we use JSON for clarity. Each document is represented as a tree $T = (V, E)$, where $V$ contains nodes (keys, values, structural elements) and $E$ encodes parent-child relationships.

\textbf{Application: LLM Consistency Evaluation.} We apply STED to evaluate consistency of LLM-generated structured outputs. Let $\mathcal{O} = \{o_1, o_2, ..., o_n\}$ denote outputs generated by an LLM for identical inputs. We compute three complementary statistics (Section~\ref{sec:stability-score}): (a) \textit{Consistency} $C_{\text{mean}}$---average STED similarity across valid pairs; (b) \textit{Stability Score} $S_\alpha$---transformed measure where higher = more stable outputs; (c) \textit{Validity} $r_v$---fraction of parseable outputs.

Tables~\ref{tab:sharegpt-consistency} and~\ref{tab:toucan-stability} report stability scores and validity across temperatures. Full $C_{\text{mean}}$ results appear in Appendix~\ref{sec:results}. The core challenge is distinguishing benign variations from critical differences that affect functionality.

\subsection{Tree Representation of Structured Data}

We transform structured documents into trees $T = (V, E)$ where each node $v \in V$ contains: \textbf{Type}: $\{\text{object}, \text{array}, \text{string}, \text{number}, \text{boolean}, \text{null}\}$; \textbf{Label}: Key name for object properties; \textbf{Value}: Raw value for primitives, child list for arrays; \textbf{Path}: Hierarchical path for unique identification.

The transformation rules are: (1) \textbf{Objects} become internal nodes with labeled edges to child values, (2) \textbf{Arrays} remain as nodes containing element lists, recursively transformed if complex, and (3) \textbf{Primitives} map to leaf nodes with actual values. This preserves array structure while enabling order-invariant matching and maintaining the semantic distinction between arrays and objects.

\subsection{Semantic Tree Edit Distance}

Our STED algorithm extends classical tree edit distance with semantic awareness through three operations on nodes $v$: \textbf{Insert} (cost $\gamma_{ins}(v)$), \textbf{Delete} (cost $\gamma_{del}(v)$), and \textbf{Update} (cost $\gamma_{upd}(v_1, v_2)$). The node similarity combines key (label) and content dimensions:
\begin{equation}
s_{node}(v_1, v_2) = w \cdot s_{key}(v_1, v_2) + (1-w) \cdot s_{children}(v_1, v_2)
\end{equation}
where $w \in [0,1]$ is the key weight (default 0.5), $s_{key}$ measures key/label similarity using embedding-based comparison, and $s_{children}$ is the normalized similarity of recursively matched children (or $s_{value}$ for primitives). The update cost is $\gamma_{upd} = 1 - s_{node}$. For objects, we employ Hungarian algorithm for optimal child matching, treating key order as semantically irrelevant. For arrays, we similarly use order-invariant matching---while JSON arrays are technically ordered, many LLM outputs exhibit semantically equivalent reorderings that should not penalize consistency scores.

\textbf{Cost Bounds.} All costs are explicitly bounded in $[0, 1]$: $\gamma_{ins}(v), \gamma_{del}(v) \in [0, 1]$ (base cost 1.0, optionally reduced by depth-decay weighting $\lambda^d$ where $\lambda \leq 1$, default $\lambda = 1$); similarity functions $s_{key}, s_{children} \in [0,1]$, so $\gamma_{upd} = 1 - s_{node} \in [0, 1]$. The normalization in Eq.~(3) divides by $\max(|C_1|, |C_2|)$, which bounds the per-level score since each matched pair contributes at most 1.0 and each unmatched node contributes at most 1.0 (insert/delete cost). Thus $d_{\text{total}} \leq \max(|C_1|, |C_2|)$, ensuring $\text{STED} \in [0,1]$.

\begin{proposition}[STED Similarity Properties]
\label{thm:similarity}
For any JSON trees $T_1, T_2$, the STED similarity $s_{\text{STED}}(T_1, T_2)$ satisfies: (i) \textbf{Boundedness}: $s_{\text{STED}}(T_1, T_2) \in [0, 1]$; (ii) \textbf{Identity}: $s_{\text{STED}}(T_1, T_2) = 1 \Rightarrow T_1 \equiv_c T_2$ (structurally identical after canonicalization). The converse holds by construction: we enforce $\text{sim}(s_1, s_2) = 1$ \emph{only} for exact string equality after canonicalization; embedding similarity is capped at $1 - \epsilon$ for non-identical strings; (iii) \textbf{Symmetry}: $s_{\text{STED}}(T_1, T_2) = s_{\text{STED}}(T_2, T_1)$; (iv) \textbf{Order Invariance}: permuting object keys does not change similarity; for arrays, we use optimal matching rather than positional comparison (a design choice for LLM consistency evaluation where element reorderings are often semantically benign).
\end{proposition}

Properties (i)--(iii) follow from cost function definitions; (iv) follows from minimum-cost bipartite matching on child multisets (Appendix~\ref{sec:similarity-proof}).

\textbf{Remark (Not a metric).} The induced distance $d = 1 - s$ satisfies non-negativity, identity, and symmetry, but triangle inequality is \textit{not} guaranteed when update costs depend on embedding similarities. We therefore treat STED as a task-motivated similarity measure, not a mathematical metric.

\subsection{Semantic Similarity Computation}

Field names are normalized (e.g., \texttt{"userName"} $\rightarrow$ \texttt{"user name"}) to capture semantic relationships. We compute similarity using embeddings with cosine similarity for texts $<300$ characters; longer texts are chunked recursively with 50-character overlaps. We used Amazon Titan Text Embeddings v2 in our implementation, though the choice of embedding model has minimal impact on STED's performance since we primarily match short field names and values where most modern embedding models achieve similar semantic understanding. The framework remains model-agnostic and can use any embedding model (e.g., all-MiniLM-L6-v2, Sentence-BERT, OpenAI embeddings).

Type-aware comparison ensures appropriate metrics: semantic for strings, exact for numbers, order-invariant for arrays. This recognizes functional equivalence (e.g., \{\texttt{"user\_name"}: \texttt{"John"}\} $\equiv$ \{\texttt{"userName"}: \texttt{"John"}\}) while detecting type violations that break compatibility. Unlike purely neural metrics requiring human validation, STED's deterministic rule-based design ensures interpretable behavior validated through controlled synthetic experiments (Section~\ref{sec:method-verification}).

\subsection{Level-wise Child Matching}

Traditional tree edit distance algorithms process nodes sequentially, which can miss better alignments when sibling order is semantically irrelevant. We address this through Hungarian algorithm-based matching at each tree level, providing optimal child-to-child assignment within each node. Note that while the classical unordered tree edit distance problem is NP-hard~\cite{Zhang1992}, our recursive level-wise approach achieves polynomial time complexity by solving optimal assignment independently at each node.

For trees $T_1$ and $T_2$ with children sets $C_1$ and $C_2$, we formulate optimal matching as an assignment problem:
\begin{multline}
\text{OptimalCost}(T_1, T_2) = \min_{\pi} \Big[ \sum_{(i,j) \in \pi} d(c_1^i, c_2^j) \\
+ \sum_{i \notin \pi} \gamma_{del}(c_1^i) + \sum_{j \notin \pi} \gamma_{ins}(c_2^j) \Big]
\end{multline}
where $\pi$ represents the matching between children, $d(\cdot, \cdot)$ is the recursive distance, and unmatched nodes incur insertion/deletion costs.

The algorithm constructs a cost matrix $M \in \mathbb{R}^{\max(|C_1|, |C_2|) \times \max(|C_1|, |C_2|)}$ where $M_{ij}$ represents the cost of matching child $i$ from $T_1$ with child $j$ from $T_2$. The Hungarian algorithm finds the minimum-cost assignment in $O(n^3)$ time, ensuring optimal child alignment at each level rather than greedy decisions.

\begin{theorem}[Level-wise Assignment Optimality]
\label{thm:hungarian}
Given trees $T_1, T_2$ with root children sets $C_1, C_2$, the Hungarian algorithm finds the optimal assignment $\pi^*$ minimizing the total matching cost \emph{at that level} in $O(\max(|C_1|, |C_2|)^3)$ time.
\end{theorem}

This provides local optimality at each tree level. The proof constructs a padded cost matrix reducing child matching to the classical assignment problem solvable by Kuhn-Munkres (Appendix~\ref{sec:proofs}). While this recursive level-wise approach does not guarantee the globally optimal unordered tree edit distance (an NP-hard problem), it provides a principled polynomial-time similarity measure that empirically correlates strongly with human judgments of structural equivalence.

\subsection{Similarity Score Normalization}

The STED algorithm normalizes the computed distance at each tree level to produce an interpretable similarity score in [0,1]:
\begin{equation}
\text{STED}(T_1, T_2) = 1 - \min\left(1, \frac{d_{\text{total}}}{\max(|C_1|, |C_2|)}\right)
\end{equation}
where $d_{\text{total}}$ is the total cost from Hungarian assignment (including insert/delete costs for unmatched nodes). This per-level normalization ensures consistent scoring regardless of tree depth.

\subsection{Consistency Aggregation}
\label{sec:stability-score}

Given $n$ generations $\mathcal{O} = \{o_1, \ldots, o_n\}$ from identical inputs, let $V \subseteq \{1,\ldots,n\}$ be indices of valid (parseable) outputs with $m = |V|$ (the number of valid outputs), and let $s_{ij} = s_{\text{STED}}(o_i, o_j)$ denote pairwise similarity. We define three complementary statistics:

\textbf{Consistency ($C_{\text{mean}}$).} The average STED similarity across valid pairs:
\begin{equation}
C_{\text{mean}}(\mathcal{O}) = \frac{2}{m(m-1)} \sum_{i<j,\, i,j \in V} s_{ij}
\end{equation}
This answers: \textit{if we sample two runs at random, how similar are they on average?}

\textbf{Validity ($r_v$).} The fraction of successful parses: $r_v(\mathcal{O}) = m/n$.

\textbf{Stability Score ($S_{\alpha}$).} For model benchmarking, we transform output variability into an intuitive score where higher = more stable:
\begin{equation}
S_{\alpha}(\mathcal{O}) = \left(\frac{1}{1 + 2\hat{\sigma}}\right)^{\alpha}, \quad \hat{\sigma} = \frac{\sigma}{\sigma_{\max}}, \quad \sigma = \text{std}(\{s_{ij}\})
\end{equation}
where $\sigma_{\max}$ is the maximum achievable std for $\binom{m}{2}$ values in $[0,1]$, and $\alpha=20$ (default). This amplifies small variations for efficient model ranking. Ablation in Appendix~\ref{sec:stability-analysis}.

\begin{proposition}[Convergence Rate]
\label{prop:convergence}
Let $\{o_i\}_{i=1}^n$ be i.i.d. samples from an LLM's output distribution. As $n \to \infty$, $C_{\text{mean}}(\mathcal{O}_n) \xrightarrow{p} C^*$ with rate $O(n^{-1/2})$, following from U-statistic theory applied to pairwise similarities.
\end{proposition}

\subsection{Computational Complexity}

For trees with $N$ nodes and maximum branching factor $B$, STED runs in $O(N \cdot B^3)$ time (Hungarian algorithm at each node) and $O(N \cdot d)$ space (embedding cache). Our optimized implementation achieves 7--22$\times$ speedup over naive recursion through memoization and embedding precomputation, maintaining 0.96 correlation with exact computation. Details in Appendix~\ref{sec:optimizations}; proofs in Appendix~\ref{sec:proofs}.

\section{Experiments and Results}

Our experimental evaluation comprises two complementary studies. First, we validate our method's effectiveness using synthetic datasets with controlled variations to verify its ability to accurately quantify similarity degradation. Second, we deploy our framework to benchmark the structured output consistency of \textbf{18 LLMs across 10 providers} on two tasks: (1) \textbf{ShareGPT Structured Output} (75 samples, 18 models) and (2) \textbf{Toucan Tool Calling} (1,006 samples, 18 models).

Models span diverse architectures: Claude family (3.5-Sonnet, 3.7-Sonnet, Haiku-4.5, Opus-4/4.5, Sonnet-4/4.5, 3.5-Haiku), Qwen (235B-A22B, 32B), Meta (Llama-3.3-70B), AWS (Nova-2-Lite), xAI (Grok-4.1-Fast), Minimax (M2), Moonshot (Mimo-V2-Flash), OpenAI (GPT-OSS-120B, GPT-4.1-Mini), and Google (Gemini-2.5-Flash-Lite). This comprehensive coverage reveals consistency patterns across model scales (32B--235B parameters), optimization strategies (efficiency vs performance), and access tiers (commercial vs free). The framework itself is model-agnostic and can evaluate any LLM with structured output capabilities.

\subsection{Experimental Setup}

\subsubsection{Validation Dataset}

We verify method effectiveness using \textbf{2,400 synthetic test cases} generated from 75 ShareGPT JSON samples with controlled modifications at 10 intensity levels (0.1--1.0): field name variants (semantic equivalents), expression variants (paraphrases), semantic variants (content changes), and structural variants (flattening/nesting). This controlled setup isolates specific variation types to validate STED's discrimination ability. Details in Appendix~\ref{sec:dataset-distribution}.

\subsubsection{Baseline Methods}

We compare STED against three established approaches: \textbf{TED}: Tree Edit Distance using Zhang-Shasha algorithm with default cost functions; \textbf{BERTScore}: Adapted for JSON by serializing structures with sorted keys, computing F1 scores from BERT embedding alignments; \textbf{DeepDiff}: Rule-based structural comparison using Deep Distance metric, converted to similarity as $1 - \text{DeepDistance}$.

\subsubsection{LLM Benchmarking Protocol}

We design a comprehensive framework to evaluate LLM consistency across two complementary tasks, summarized in Table~\ref{tab:experimental-protocol}.

\begin{table}[h]
\centering
\caption{Experimental Protocol Summary}
\label{tab:experimental-protocol}
\scriptsize
\begin{tabular}{@{}lcc@{}}
\toprule
\textbf{Parameter} & \textbf{ShareGPT} & \textbf{Toucan} \\
\midrule
Samples & 75 & 1,006 \\
Temperatures & 11 (0.0--1.0) & 11 (0.0--1.0) \\
Runs per temp & 10 & 10 \\
\textbf{Outputs/model} & \textbf{8,250} & \textbf{110,660} \\
Models evaluated & 18 & 18 \\
Unique tools & -- & 993 \\
\bottomrule
\end{tabular}

\vspace{0.3em}
\raggedright\scriptsize Total: 2.2M+ outputs across all models.
\end{table}

\textbf{Model Coverage.} We evaluate 18 models across 10 providers. Main paper tables report 12 representative models for readability; Appendix tables include all 18 models. All stability scores ($S_\alpha$) are computed from hybrid STED similarities ($w=0.5$) over valid output pairs.

\textbf{ShareGPT Structured Output}: 75 JSON generation samples (5 excluded due to malformed ground truth), each executed 10 times across 11 temperatures (0.0--1.0), yielding \textbf{8,250 outputs per model}. \textbf{Toucan Tool Calling}: 1,006 samples from the Toucan benchmark~\cite{toucan2024} (993 unique tools, 56\% multilingual) with the same protocol, yielding \textbf{110,660 outputs per model}. Tool calling complements free-form JSON as a constrained task where format is SDK-enforced but models vary in tool selection and parameters. Combined across all models, our evaluation comprises over \textbf{2.2 million LLM outputs}. Dataset details in Appendix~\ref{sec:dataset-distribution}.

\textbf{Validity Definition.} We define \textit{validity} ($r_v$) as successful parse rate: (1) \textbf{ShareGPT}: output must be valid JSON (parseable by standard JSON decoder after markdown fence extraction); (2) \textbf{Toucan}: tool call must have valid structure (tool name present, arguments parseable as JSON). Note: validity measures \textit{syntactic correctness}, not schema compliance or semantic accuracy. Low validity for some models on Toucan (e.g., Gemini 54\%, Llama 33\%) reflects API/SDK differences in tool-calling format enforcement across providers, not JSON parsing failures per se.

We employ three evaluation modes: (1) structural consistency measuring format adherence, (2) semantic consistency evaluating content preservation, and (3) hybrid approach with equal weighting ($w = 0.5$). This multi-faceted evaluation reveals how temperature variations affect different aspects of structured output generation.

\begin{figure}[t]
\centering
\includegraphics[width=0.95\columnwidth]{figures/schema_variation_analysis.png}
\caption{Similarity scores across schema variation types for four consistency evaluation methods. Field name changes (0.1-1.0 ratios) show gradual similarity degradation, with BERTScore maintaining highest scores. For structural variations, STED achieves zero similarity on nested changes but moderate scores on flat structures. All methods use similarity scores ranging from 0 to 1.}
\label{fig:schema_variation_analysis}
\end{figure}

\begin{figure}[t]
\centering
\includegraphics[width=0.95\columnwidth]{figures/similarity_progression_combined.png}
\caption{Similarity score progression under semantic variations (left) and expression variations (right). Ideal behavior: sensitivity to semantic changes, robustness to expression changes. STED achieves optimal balance with controlled degradation for semantic variations while maintaining high scores for expression variations, aligning with human consistency perception. (TED remains constant at 1.0 as it measures only structural similarity, not content.)}
\label{fig:similarity_progression_combined}
\end{figure}

\subsection{Method Validation Results}
\label{sec:method-verification}

\subsubsection{Schema Variation Analysis}
Figure~\ref{fig:schema_variation_analysis} reveals how methods handle structural variations in LLM outputs. For field name variations using semantic equivalents, STED maintains stable similarity across all variation ratios, correctly recognizing functional equivalence. TED degrades significantly due to lack of semantic understanding, while BERTScore's excessive permissiveness may mask important structural differences. For structural modifications (flattening and nesting changes), STED correctly assigns zero similarity, recognizing these as breaking changes for downstream systems. Other methods problematically tolerate these changes with non-zero scores that could allow incompatible outputs to pass consistency checks.

\subsubsection{Content Variation Analysis}
Figure~\ref{fig:similarity_progression_combined} evaluates content variation handling, critical for applications like content moderation where semantic accuracy impacts outcomes. For semantic variations, STED (0.954$\pm$0.039) and BERTScore (0.958$\pm$0.025) show statistically equivalent calibrated sensitivity ($p = 0.600$), while DeepDiff over-reacts (0.799$\pm$0.165) and TED remains insensitive (1.0, structure-only). For expression variations, STED achieves superior robustness (0.981$\pm$0.017, $p < 0.001$), correctly preserving high similarity for paraphrases, unlike DeepDiff which incorrectly penalizes valid rewording (0.805$\pm$0.164).\footnote{Detailed statistical analysis in Appendix~\ref{sec:statistical-analysis}.}

STED's differential response---5.2\% degradation for semantic changes versus 1.9\% for expression changes---enables threshold-based detection of genuine semantic drift while accepting natural language variation, providing human-aligned consistency assessment for practical applications.

\subsubsection{Human Validation Study}
\label{sec:human-validation}

To validate STED on real LLM outputs, we conducted a best-match selection study with three annotators. Of 500 Toucan samples processed, methods agreed on the best match in 344 cases (68.8\%); we evaluated 74 disagreement cases where method choice matters (agreement cases provide no signal for comparison). Three annotators (Fleiss' $\kappa = 0.759$, substantial agreement) independently selected the most similar candidate to ground truth, with methods anonymized and order randomized. Using majority voting across annotators, STED achieves highest alignment with human judgment: \textbf{87.8\%} win rate (65/74), followed by BERTScore (74.3\%), DeepDiff (47.3\%), and TED (27.0\%). On the 59 items where all three annotators unanimously agreed, STED achieves \textbf{91.5\%} alignment. STED significantly outperforms DeepDiff and TED ($p < 0.001$, McNemar's test); the STED-BERTScore gap (13.5pp) approaches statistical significance ($p = 0.087$). This complements synthetic validation: controlled experiments test discrimination (2,400 cases), while human validation confirms alignment on real outputs.

\subsection{LLM Consistency Results}

Our comprehensive evaluation across 18 models on two distinct tasks reveals critical insights into consistency patterns, providing guidance for model selection and parameter tuning.

\subsubsection{ShareGPT Structured Output}
Table~\ref{tab:sharegpt-consistency} shows stability scores across temperatures for 12 models. Key findings: (1) \textbf{stability degrades with temperature}---Claude-3.5-Haiku drops from 0.80 to 0.24; (2) Claude-3.5-Sonnet/3.7-Sonnet show best temperature robustness ($\Delta \approx -0.12$).

\textbf{Prompt Characteristic Analysis.} We analyzed 75 ShareGPT prompts to identify factors associated with output inconsistency, aggregating stability scores across all temperatures and 18 models. Table~\ref{tab:ambiguity-main} compares top and bottom quartiles by stability score. The most notable difference is in question mark frequency: inconsistent samples average 2.8$\times$ more question marks, suggesting prompts with multiple clarification requests may introduce ambiguity. Prompt length and prose-heavy text show no meaningful difference. Correlation analysis confirms no significant relationship between prompt characteristics and consistency ($r = 0.03$, $p = 0.80$), indicating consistency is primarily driven by model capabilities and task difficulty rather than prompt characteristics.

\begin{table}[h]
\centering
\caption{Prompt Factors vs Consistency (ShareGPT, All Temperatures)}
\label{tab:ambiguity-main}
\scriptsize
\begin{tabular}{@{}lccc@{}}
\toprule
\textbf{Factor} & \textbf{Consistent} & \textbf{Inconsist.} & \textbf{Ratio} \\
\midrule
Question marks (avg) & 0.8 & 2.4 & 2.8$\times$ \\
Prompt length (chars) & 15,298 & 15,877 & 1.0$\times$ \\
Prose-heavy$^\dagger$ & 5.0\% & 5.0\% & 1.0$\times$ \\
\bottomrule
\end{tabular}

\vspace{0.3em}
\raggedright\scriptsize $^\dagger$Avg sentence $>$20 words, no structure markers. Consistent: $S_\alpha \geq 0.51$ (N=20); Inconsistent: $S_\alpha \leq 0.35$ (N=20). Quartile-based thresholds across 18 models, all temperatures.
\end{table}

\subsubsection{Toucan Tool Calling}
Table~\ref{tab:toucan-stability} presents stability scores across temperatures for 12 models. Claude-Opus-4 achieves highest stability (0.71 at T=0.0), while Nova-2-Lite/Claude-3.7-Sonnet show best temperature robustness ($\Delta \approx -0.04$). Qwen3-235B-A22B leads validity (98.1\%).

\textbf{Consistency Diagnostics: Parameter Complexity as Primary Predictor.} A key contribution of this work is identifying \textit{what causes} inconsistency. Analyzing 1,006 tool calling samples across 18 models and all temperatures, we find that \textbf{structural complexity of parameters strongly predicts inconsistency}. Table~\ref{tab:complexity-main} shows over-representation of complex types in inconsistent samples: lists are \textbf{27$\times$} more likely to appear in inconsistent outputs, and samples with any complex type are 29$\times$ more likely to be inconsistent. (A finer-grained analysis at T=0.5 in Appendix Table~\ref{tab:complexity-correlation} shows 60$\times$ for deeply nested structures specifically.)

\begin{table}[h]
\centering
\caption{Parameter Complexity vs Inconsistency (Toucan, All Temperatures)}
\label{tab:complexity-main}
\scriptsize
\begin{tabular}{@{}lccc@{}}
\toprule
\textbf{Parameter Type} & \textbf{Consistent} & \textbf{Inconsist.} & \textbf{Ratio} \\
\midrule
Has List Params & 0.4\% & 10.7\% & 27$\times$ \\
Has Any Complex Type & 0.4\% & 11.5\% & 29$\times$ \\
Has Nested Structure & 0.0\% & 3.6\% & $>$3.6\% \\
Has Dict Params & 0.0\% & 2.4\% & $>$2.4\% \\
\midrule
Mean Max Depth & 0.80 & 0.99 & 1.2$\times$ \\
\bottomrule
\end{tabular}

\vspace{0.3em}
\raggedright\scriptsize Consistent: $S_\alpha \geq 0.85$ (N=252); Inconsistent: $S_\alpha \leq 0.58$ (N=252). Quartile-based thresholds across 18 models, all temperatures.
\end{table}

This finding enables actionable diagnostics: (1) it \textbf{justifies STED's structure-aware approach}---simple exact matching would miss semantic equivalences in nested structures where models vary; (2) it provides \textbf{intervention guidance}---tools with complex parameter schemas require schema simplification, more capable models, or lower temperatures; (3) it enables \textbf{proactive risk assessment}---practitioners can predict consistency challenges from schema inspection alone. This ratio remains consistent (22--31$\times$) across all temperatures, confirming parameter complexity is an intrinsic difficulty factor independent of sampling variability.

\textbf{Task-Specific Model Differences.} Models exhibit dramatically different consistency across tasks. Llama-3.3-70B achieves 98.4\% validity on ShareGPT structured output but only 32.7\% on Toucan tool calling---a 3$\times$ gap. This underscores the importance of task-specific evaluation rather than relying on general capability benchmarks.

\begin{table}[t]
\centering
\caption{ShareGPT Structured Output: Stability Score ($S_\alpha$) by Temperature (75 samples, 12 models)}
\label{tab:sharegpt-consistency}
\resizebox{\columnwidth}{!}{%
\begin{tabular}{@{}lccc|c|c@{}}
\toprule
\textbf{Model} & \textbf{T=0.0} & \textbf{T=0.5} & \textbf{T=1.0} & \textbf{$\Delta$} & \textbf{Valid.} \\
\midrule
Claude-3.5-Haiku & \textbf{.80} & .34 & .24 & $-$.57 & 99.5\% \\
Gemini-2.5-Flash-Lite & .78 & .30 & .25 & $-$.53 & \textbf{100\%} \\
Claude-Haiku-4.5 & .75 & .37 & .30 & $-$.44 & \textbf{100\%} \\
Claude-Opus-4 & .74 & .43 & .33 & $-$.40 & 97.6\% \\
Claude-Sonnet-4.5 & .64 & .38 & .34 & $-$.30 & 89.6\% \\
Claude-Opus-4.5 & .63 & .45 & .41 & \textbf{$-$.22} & 96.1\% \\
GPT-4.1-Mini & .50 & .32 & .26 & $-$.24 & \textbf{100\%} \\
Qwen3-32B & .49 & .29 & .20 & $-$.30 & 98.6\% \\
Claude-3.5-Sonnet & .47 & .41 & .33 & \textbf{$-$.13} & \textbf{100\%} \\
Claude-3.7-Sonnet & .47 & .39 & .36 & \textbf{$-$.12} & \textbf{100\%} \\
Qwen3-235B-A22B & .44 & .36 & .31 & \textbf{$-$.13} & \textbf{100\%} \\
Claude-Sonnet-4 & .40 & .29 & .27 & \textbf{$-$.13} & 73.1\% \\
\bottomrule
\end{tabular}%
}

\vspace{0.3em}
\raggedright\scriptsize Higher = more stable. $\Delta$ = change from T=0.0 to T=1.0. Full temperature breakdown in Appendix Table~\ref{tab:overall-consistency-temp}.
\end{table}

\begin{table}[t]
\centering
\caption{Toucan Tool Calling: Stability Score ($S_\alpha$) by Temperature (1,006 samples, 12 models)}
\label{tab:toucan-stability}
\resizebox{\columnwidth}{!}{%
\begin{tabular}{@{}lccc|c|c@{}}
\toprule
\textbf{Model} & \textbf{T=0.0} & \textbf{T=0.5} & \textbf{T=1.0} & \textbf{$\Delta$} & \textbf{Valid.} \\
\midrule
Claude-Opus-4 & \textbf{.71} & \textbf{.63} & \textbf{.59} & $-$.11 & 97.1\% \\
Claude-Sonnet-4.5 & .68 & .59 & .54 & $-$.14 & 96.8\% \\
Claude-Sonnet-4 & .68 & .63 & .58 & $-$.10 & 97.0\% \\
Claude-Opus-4.5 & .66 & .58 & .55 & $-$.11 & 96.7\% \\
Claude-Haiku-4.5 & .65 & .51 & .45 & $-$.20 & 95.3\% \\
Nova-2-Lite & .63 & .62 & .60 & \textbf{$-$.04} & 96.5\% \\
Mimo-V2-Flash & .58 & .51 & .41 & $-$.17 & 97.7\% \\
Claude-3.5-Haiku & .58 & .51 & .49 & $-$.09 & 97.7\% \\
Qwen3-235B-A22B & .56 & .54 & .48 & $-$.09 & \textbf{98.1\%} \\
Claude-3.5-Sonnet & .53 & .50 & .49 & \textbf{$-$.05} & 97.4\% \\
Claude-3.7-Sonnet & .52 & .50 & .49 & \textbf{$-$.04} & 96.7\% \\
Grok-4.1-Fast & .45 & .43 & .43 & \textbf{$-$.03} & 94.8\% \\
\bottomrule
\end{tabular}%
}

\vspace{0.3em}
\raggedright\scriptsize Higher = more stable. Claude-Opus-4 most stable; Nova-2-Lite/Claude-3.7-Sonnet most temperature-robust ($\Delta \approx -0.04$). Full $C_\text{mean}$ results in Appendix Table~\ref{tab:toucan-full-temp}.
\end{table}

\textbf{Interpreting $S_\alpha$ vs $C_\text{mean}$.} The $S_\alpha$ values amplify variance for model ranking; underlying $C_\text{mean}$ is higher (e.g., $S_\alpha$=0.71 $\rightarrow$ $C_\text{mean}$$\approx$0.95). Small $C_\text{mean}$ gaps (0.95 vs 0.92) yield meaningful $S_\alpha$ differences (0.71 vs 0.52), enabling discrimination. Full $C_\text{mean}$ in Appendix.

\subsection{Ablation Studies}

\subsubsection{Steepness Parameter ($\alpha$)}
\label{sec:ablation-steepness}

Table~\ref{tab:alpha-ablation-main} shows the effect of $\alpha$ on score discrimination. We select $\alpha = 20$ as default, achieving full $[0,1]$ range normalization while retaining 95\% of peak discrimination. Model rankings are robust to $\alpha$ choice: Spearman $\rho = 0.998$ on Toucan and $\rho = 0.987$ on ShareGPT across $\alpha \in \{5, 10, 15, 20, 25, 30, 40, 50\}$. Detailed analysis in Appendix~\ref{sec:stability-analysis}.

\begin{table}[h]
\centering
\caption{Effect of $\alpha$ on discrimination}
\label{tab:alpha-ablation-main}
\scriptsize
\begin{tabular}{@{}lccccc@{}}
\toprule
\textbf{Metric} & $\alpha=5$ & $\alpha=10$ & $\alpha=15$ & $\alpha=20$ & $\alpha=25$ \\
\midrule
$S_\alpha(\hat{\sigma}=0.01)$ & 0.905 & 0.819 & 0.741 & 0.670 & 0.607 \\
$S_\alpha(\hat{\sigma}=0.10)$ & 0.402 & 0.162 & 0.065 & 0.026 & 0.011 \\
Discrimination & 0.503 & 0.657 & \textbf{0.676} & 0.644 & 0.596 \\
\bottomrule
\end{tabular}
\end{table}

\subsubsection{Embedding Model}
\label{sec:ablation-embedding}

We evaluate STED with nine embedding models spanning three architectures on 200 synthetic samples: Sentence-Transformers (MiniLM-L6, MPNet), BGE (base, large), E5 (base, large), and Amazon Titan (256d, 512d, 1024d). Table~\ref{tab:embedding-ablation} shows extremely high correlation across all pairs ($r > 0.999$, $\rho > 0.995$), confirming that embedding model choice has minimal impact on STED's consistency measurements.

\begin{table}[h]
\centering
\caption{Embedding model ablation: pairwise Pearson correlation ($r$)}
\label{tab:embedding-ablation}
\scriptsize
\begin{tabular}{@{}lccccc@{}}
\toprule
& \textbf{MiniLM} & \textbf{MPNet} & \textbf{BGE-L} & \textbf{E5-L} & \textbf{Titan-1024} \\
\midrule
MiniLM-L6 & -- & 1.000 & 1.000 & 0.999 & 1.000 \\
MPNet & & -- & 1.000 & 0.999 & 1.000 \\
BGE-large & & & -- & 1.000 & 1.000 \\
E5-large & & & & -- & 0.999 \\
\midrule
Min $\rho$ & \multicolumn{5}{c}{0.995 (E5 vs Titan)} \\
\bottomrule
\end{tabular}

\vspace{0.3em}
\raggedright\scriptsize All 36 pairwise correlations exceed $r=0.999$ and $\rho=0.995$ across 9 embedding models (200 samples). Full results include BGE-base, E5-base, Titan-256d, Titan-512d.
\end{table}

\subsection{Practitioner Guidelines}
\label{sec:practitioner-guidelines}

We synthesize our findings into actionable deployment recommendations.

\textbf{Model \& Temperature Selection.} For maximum consistency: Claude-Opus-4 at T=0.0--0.1 ($S_\alpha$=0.71). For temperature robustness: Claude-3.7-Sonnet ($\Delta$=$-$0.04 across T range). For maximum validity: Qwen3-235B-A22B (98--100\% parse success). Cost-effective: Claude-Haiku-4.5 at T=0.2 ($S_\alpha$=0.76).

\textbf{Schema Design.} Parameter complexity predicts inconsistency (27--29$\times$ for complex types): (1) \textit{flatten nested structures} to key-value pairs; (2) \textit{avoid list parameters}---use fixed fields or pagination; (3) \textit{prefer primitives} (string/number/boolean); (4) \textit{include examples} in tool descriptions.

\textbf{Production Monitoring.} Establish $C_\text{mean}$ baseline over 50--100 outputs; alert on $>$5\% deviation. Track failure categories (string drift, tool count, tool selection, complex params) to diagnose degradation. Apply stricter thresholds for complex-parameter tools ($S_\alpha > 0.6$) vs simple tools ($S_\alpha > 0.4$).

\textbf{When to Use STED.} STED provides value when outputs contain semantic variations in field names/values, array ordering varies, or you need to distinguish benign variations from breaking changes. For exact-match requirements, simple equality suffices.

\section{Conclusion}

We presented STED (Semantic Tree Edit Distance), a similarity metric for complex JSON structures that combines structural tree comparison with semantic understanding via embeddings and order-invariant matching. Human evaluation (87.8\% alignment via majority vote, 3 annotators, Fleiss' $\kappa=0.759$) confirms STED captures human intuition about structured data similarity, significantly outperforming DeepDiff and TED ($p<0.001$).

Applying STED to LLM consistency evaluation across 18 models (2.2M+ outputs), we introduce a \textbf{consistency diagnostics framework} identifying \textit{why} outputs are inconsistent: parameter complexity is strongly associated with inconsistency (27--29$\times$ higher for complex types), with distinct failure patterns (string drift 39\%, tool count 35\%, tool selection 16\%, complex parameters 6\%) enabling targeted interventions. Task-specific differences are dramatic (Llama-3.3-70B: 99\% vs 33\% validity across tasks).

While our validation focuses on LLM outputs, STED's design principles---semantic matching, order invariance, and hierarchical comparison---address fundamental challenges in JSON similarity that arise in other domains (API testing, configuration management). Empirical validation in these broader contexts remains future work.

\textbf{Limitations.} We focus on JSON format; while the algorithm generalizes to other hierarchical formats (XML, HTML), empirical validation remains future work. The human evaluation targets discriminative cases where methods disagree, complementing synthetic validation on representative samples. Embedding-based semantic matching may miss domain-specific equivalences without fine-tuning.

\bibliography{neurips_2025_camera_ready}
\bibliographystyle{icml2025}

%%%%%%%%%%%%%%%%%%%%%%%%%%%%%%%%%%%%%%%%%%%%%%%%%%%%%%%%%%%%
\appendix
\section{Theoretical Proofs}
\label{sec:proofs}

\subsection{Proof of Similarity Properties (Proposition~\ref{thm:similarity})}
\label{sec:similarity-proof}

\textbf{(i) Boundedness:} We prove $s_{\text{STED}} \in [0,1]$ by showing that cost functions are bounded and normalization is correct.

\textit{Cost bounds:} (a) Similarity functions: cosine similarity is clipped to $[0,1]$, so $\gamma_{struct} = 1 - \text{sim}_{struct} \in [0,1]$ and $\gamma_{content} = 1 - \text{sim}_{content} \in [0,1]$. Thus $\gamma_{upd} = w \cdot \gamma_{struct} + (1-w) \cdot \gamma_{content} \in [0,1]$. (b) Insert/delete: $\gamma_{ins}, \gamma_{del}$ are explicitly clamped to $[0,1]$ (base cost 1.0, optionally reduced by depth-decay $\lambda^d$ where $\lambda \leq 1$).

\textit{Normalization correctness:} At each tree level with children sets $C_1, C_2$ of sizes $n_1, n_2$, the Hungarian assignment matches $\min(n_1, n_2)$ pairs (each contributing $\leq 1$) plus $|n_1 - n_2|$ unmatched nodes (each contributing insert/delete cost $\leq 1$). Thus $d_{\text{total}} \leq \max(n_1, n_2)$, and division by $\max(|C_1|, |C_2|)$ in Eq.~(3) yields $\text{STED} \in [0,1]$.

\textbf{(ii) Identity:}
$(\Rightarrow)$ If $s_{\text{STED}}(T_1, T_2) = 1$, then $d_{\text{total}} = 0$ in Eq.~(3). This requires: no insertions/deletions (since $\gamma_{\text{ins}}, \gamma_{\text{del}} > 0$), and all updates have zero cost, meaning $S_{\text{struct}}(n_1, n_2) = 1$ and $S_{\text{value}}(n_1, n_2) = 1$ for all matched pairs. Since our implementation returns $\text{sim}(s_1, s_2) = 1$ only for exact string equality (after canonicalization), this implies all labels and values are identical after normalization, hence $T_1 \equiv_c T_2$.

$(\Leftarrow)$ If $T_1 \equiv_c T_2$ (identical after canonicalization), then the identity mapping yields $d_{\text{total}} = 0$: all nodes match exactly with $S_{\text{struct}} = S_{\text{value}} = 1$, so $s_{\text{STED}} = 1$.

\textbf{Remark (Implementation guarantee):} To ensure the identity property holds exactly, our implementation enforces: (1) $\text{sim}(s_1, s_2) = 1$ \emph{only} when $s_1 = s_2$ after canonicalization (case normalization, whitespace trimming); (2) for non-identical strings, embedding-based cosine similarity is capped at $1 - \epsilon$ (we use $\epsilon = 10^{-6}$) to prevent degenerate embedding cases from yielding spurious identity. This enables meaningful graded similarity for semantically similar strings (e.g., ``email'' vs ``email\_address'' $\approx 0.85$) while guaranteeing the mathematical identity property.

\textbf{(iii) Symmetry:} The cost functions are symmetric by construction: $\gamma_{\text{ins}}(n) = \gamma_{\text{del}}(n)$ (same formula), and similarity functions are symmetric. For any matching $\pi: T_1 \to T_2$, the inverse matching $\pi^{-1}: T_2 \to T_1$ has equal cost.

\textbf{(iv) Order Invariance:} For objects and arrays, children are matched via minimum-cost bipartite assignment (Hungarian algorithm) on the multiset of children. The assignment considers all possible pairings and selects the minimum-cost matching. Since the cost function only depends on the children themselves (not their positions), any permutation of children in either tree results in the same set of pairwise costs, hence the same optimal matching cost. Therefore, $s_{\text{STED}}$ is invariant to permutations of object keys or array elements.

\subsection{Proof of Level-wise Assignment Optimality (Theorem~\ref{thm:hungarian})}

We reduce child matching to the classical assignment problem.

\textbf{Step 1: Cost Matrix Construction.}
Construct matrix $M \in \mathbb{R}^{m \times m}$ where $m = \max(n_1, n_2)$:
\[
M[i][j] = \begin{cases}
d_{\text{STED}}(c_i^{(1)}, c_j^{(2)}) & \text{if } i \leq n_1, j \leq n_2 \\
\gamma_{\text{del}}(c_i^{(1)}) & \text{if } i \leq n_1, j > n_2 \\
\gamma_{\text{ins}}(c_j^{(2)}) & \text{if } i > n_1, j \leq n_2 \\
0 & \text{otherwise}
\end{cases}
\]

\textbf{Step 2: Equivalence.} Matching real node $i$ to dummy $j$ represents deletion (cost $\gamma_{\text{del}}$); matching dummy $i$ to real node $j$ represents insertion (cost $\gamma_{\text{ins}}$).

\textbf{Step 3: Complexity.} The Kuhn-Munkres algorithm solves $m \times m$ assignment in $O(m^3)$ time and $O(m^2)$ space, guaranteeing optimality for the child assignment at that level via complementary slackness.

\subsection{Proof of Mean Consistency Convergence (Proposition~\ref{prop:convergence})}

The mean pairwise consistency $C_{\text{mean}} = \frac{2}{m(m-1)} \sum_{i<j} s_{ij}$ is a U-statistic of order 2 with kernel $h(o_i, o_j) = s(T_i, T_j)$.

\textbf{Step 1: U-statistic structure.}
By Hoeffding's theorem for U-statistics, $\sqrt{n}(C_{\text{mean},n} - \mu^*) \xrightarrow{d} \mathcal{N}(0, 4\zeta_1^2)$ where $\mu^* = \mathbb{E}[s(T_1, T_2)]$ and $\zeta_1^2 = \text{Var}[\mathbb{E}[s(T_1, T_2) | o_1]]$.

\textbf{Step 2: Rate bound.}
By the CLT for U-statistics: $|C_{\text{mean},n} - \mu^*| = O_p(n^{-1/2})$.

This establishes that mean pairwise consistency converges to the population value at the standard parametric rate, justifying finite-sample comparisons with $n \geq 10$ generations per condition.

\subsection{STED Algorithm Pseudocode}
\label{sec:pseudocode}

Algorithm~\ref{alg:sted} provides the complete STED computation, clarifying how per-level normalization composes across the tree.

\begin{algorithm}[h]
\caption{STED: Semantic Tree Edit Distance}
\label{alg:sted}
\begin{algorithmic}[1]
\REQUIRE Trees $T_1, T_2$; key weight $w \in [0,1]$
\ENSURE Similarity $s \in [0,1]$
\STATE \textbf{function} STED($T_1$, $T_2$, $w$)
\IF{$T_1$ and $T_2$ are both primitives}
    \IF{$\text{canonicalize}(T_1) = \text{canonicalize}(T_2)$}
        \RETURN 1.0 \COMMENT{Exact match after canonicalization}
    \ELSE
        \STATE $s_{\text{sem}} \gets \min(\text{embed\_sim}(T_1, T_2), 1-\epsilon)$
        \RETURN $s_{\text{sem}}$ \COMMENT{Capped embedding similarity}
    \ENDIF
\ENDIF
\STATE $C_1 \gets \text{children}(T_1)$; $C_2 \gets \text{children}(T_2)$
\STATE $m \gets \max(|C_1|, |C_2|)$
\STATE \textbf{Build cost matrix} $M \in \mathbb{R}^{m \times m}$:
\FOR{$i = 1$ to $|C_1|$, $j = 1$ to $|C_2|$}
    \STATE $s_{ij} \gets \text{STED}(C_1[i], C_2[j], w)$ \COMMENT{Recursive}
    \STATE $M[i,j] \gets 1 - s_{ij}$ \COMMENT{Convert to cost}
\ENDFOR
\FOR{$i = 1$ to $|C_1|$, $j > |C_2|$}
    \STATE $M[i,j] \gets \gamma_{\text{del}}(C_1[i])$ \COMMENT{Deletion cost}
\ENDFOR
\FOR{$i > |C_1|$, $j = 1$ to $|C_2|$}
    \STATE $M[i,j] \gets \gamma_{\text{ins}}(C_2[j])$ \COMMENT{Insertion cost}
\ENDFOR
\STATE $\pi^* \gets \text{Hungarian}(M)$ \COMMENT{Optimal assignment}
\STATE $d_{\text{total}} \gets \sum_{(i,j) \in \pi^*} M[i,j]$
\STATE $s_{\text{children}} \gets \max(0, 1 - d_{\text{total}} / m)$ \COMMENT{Eq.~(3)}
\STATE $s_{\text{key}} \gets \text{key\_similarity}(T_1, T_2)$ \COMMENT{Label match}
\RETURN $w \cdot s_{\text{key}} + (1-w) \cdot s_{\text{children}}$
\end{algorithmic}
\end{algorithm}

\textbf{Key clarifications:} (1) \textbf{Per-level normalization} (line 21): $s_{\text{children}} = \max(0, 1 - d_{\text{total}}/m)$ which is equivalent to Eq.~(3)'s formulation $1 - \min(1, d_{\text{total}}/m)$; the $\max(0,\cdot)$ ensures non-negativity while cost bounds guarantee this is never triggered. (2) \textbf{Global composition}: The recursive structure means deep mismatches contribute through their parent's cost matrix---a mismatch at depth $d$ affects ancestor similarity scores multiplicatively through the recursion. (3) \textbf{Depth-decay} (optional): $\gamma_{\text{ins/del}}$ can include factor $\lambda^d$ to weight shallow mismatches more heavily; we use $\lambda=1$ (uniform) by default.

\section{Dataset Characteristics and Distribution Analysis}
\label{sec:dataset-distribution}

We provide a comprehensive analysis of the 75 valid ShareGPT samples used for benchmarking and synthetic validation. From 80 initial samples, 5 were excluded due to malformed ground truth JSON, leaving 75 valid samples that generated our 2,400 synthetic test cases. This analysis demonstrates the representativeness and diversity of our dataset across multiple dimensions of JSON structural complexity.

\subsection{Structural Complexity Distribution}
\label{subsec:structural-complexity}

The distribution of JSON depths across our base samples shows the majority (58.7\%) at depth 4, aligning with typical enterprise API structures. Depths range from 2 (10.7\%, simple flat structures) to 7 (1.3\%, maximum complexity). Mean depth: 4.0$\pm$1.0.

Field count distribution demonstrates coverage from simple configurations to complex documents: 1--10 fields (8.0\%), 11--25 fields (29.3\%), 26--50 fields (36.0\%, most common), 51--100 fields (21.3\%), and 100+ fields (5.3\%). Mean: 42.3$\pm$37.8 fields.

\subsection{Field Type and Complexity Analysis}
\label{subsec:field-type-analysis}

Distribution of field types across all samples: String (7,276, 68.2\%), Integer (1,968, 18.4\%), Array (907, 8.5\%), Object (524, 4.9\%).

Statistical summary of complexity metrics: Max Depth (min 2, max 7, mean 4.0$\pm$1.0), Total Fields (min 4, max 228, mean 42.3$\pm$37.8), Total Nodes (min 10, max 320, mean 64.2$\pm$44.9), Arrays (min 0, max 11, mean 4.7$\pm$2.7), Nested Objects (min 0, max 98, mean 12.1$\pm$13.6).

\subsection{ShareGPT Structured Output Dataset Statistics}
\label{subsec:sharegpt-statistics}

Table~\ref{tab:sharegpt-dataset-stats} summarizes our ShareGPT evaluation dataset. From 80 initial samples, 5 were excluded due to malformed ground truth JSON (truncated or invalid syntax), leaving \textbf{75 valid samples} drawn from two ShareGPT-formatted datasets: \textit{Quiz Generation} (48 samples, 64.0\%) containing structured quiz questions with answers, and \textit{Structured Output} (27 samples, 36.0\%) containing various JSON generation tasks.

\begin{table}[h]
\centering
\caption{ShareGPT Dataset Statistics}
\label{tab:sharegpt-dataset-stats}
\scriptsize
\begin{tabular}{@{}lrl@{}}
\toprule
\textbf{Metric} & \textbf{Value} & \textbf{Details} \\
\midrule
\multicolumn{3}{l}{\textit{Sample Composition}} \\
Total samples & 75 & (5 excluded$^\dagger$) \\
\quad Quiz Generation & 48 & 64.0\% \\
\quad Structured Output & 27 & 36.0\% \\
\midrule
\multicolumn{3}{l}{\textit{Structure Statistics}} \\
Max depth & 4.0$\pm$1.0 & range 2--7 \\
Total fields & 42.3$\pm$37.8 & range 4--228 \\
Arrays per sample & 4.7$\pm$2.7 & range 0--11 \\
Nested objects & 12.1$\pm$13.6 & range 0--98 \\
\midrule
\multicolumn{3}{l}{\textit{Field Type Distribution}} \\
\quad string & 7,276 & 68.2\% \\
\quad integer & 1,968 & 18.4\% \\
\quad array & 907 & 8.5\% \\
\quad object & 524 & 4.9\% \\
\midrule
\multicolumn{3}{l}{\textit{Input Characteristics}} \\
Prompt length & 4,330$\pm$8,963 & mean$\pm$std (tokens) \\
\quad Range & 172--73,522 & tokens \\
\quad Median & 1,795 & tokens \\
\bottomrule
\end{tabular}

\vspace{0.3em}
\raggedright\scriptsize $^\dagger$5 samples excluded due to malformed ground truth JSON (truncated data or invalid syntax).
\end{table}

\textbf{Structural Diversity.} The ShareGPT samples exhibit diverse structural complexity: depths range from 2 (simple flat structures) to 7 (deeply nested), with the majority (58.7\%) at depth 4, aligning with typical enterprise API structures. Field counts span from simple 4-field configurations to complex 228-field documents, ensuring evaluation across varied JSON complexity levels.

\textbf{Input Characteristics.} ShareGPT prompts are substantially longer than Toucan queries (median 1,795 vs 105 tokens), reflecting detailed task descriptions and context. The high variance (std 8,963 tokens) spans from concise 172-token requests to comprehensive 73,522-token documents, testing models across diverse input complexity levels.

\subsection{Toucan Tool Calling Dataset Statistics}
\label{subsec:toucan-statistics}

Table~\ref{tab:toucan-dataset-stats} summarizes our Toucan evaluation subset. The 1,006 samples are drawn from two single-turn subsets: \textit{single-turn-original} (515, 51.2\%) and \textit{single-turn-diversify} (491, 48.8\%), selected to focus on direct tool invocation without multi-turn conversation complexity.

\begin{table}[h]
\centering
\caption{Toucan Dataset Statistics}
\label{tab:toucan-dataset-stats}
\scriptsize
\begin{tabular}{@{}lrl@{}}
\toprule
\textbf{Metric} & \textbf{Value} & \textbf{Details} \\
\midrule
\multicolumn{3}{l}{\textit{Sample Composition}} \\
Total samples & 1,006 & -- \\
Unique tools & 993 & across all samples \\
Tools per sample & 95.7 & mean available tools \\
Multilingual & 56.4\% & non-ASCII queries \\
\midrule
\multicolumn{3}{l}{\textit{Tool Call Distribution}} \\
Total calls & 1,583 & across all samples \\
Calls per sample & 1.57 & mean (range 1--14) \\
\quad 1 call & 61.9\% & 623 samples \\
\quad 2 calls & 23.5\% & 236 samples \\
\quad 3 calls & 11.9\% & 120 samples \\
\quad 4+ calls & 2.7\% & 27 samples \\
\midrule
\multicolumn{3}{l}{\textit{Parameter Statistics}} \\
Total parameters & 2,366 & across all calls \\
Params per call & 1.48 & mean (range 0--19) \\
\quad string & 48.7\% & 1,152 params \\
\quad float & 24.8\% & 587 params \\
\quad int & 18.2\% & 431 params \\
\quad list & 4.3\% & 102 params \\
\quad bool & 3.1\% & 73 params \\
\quad dict & 0.9\% & 21 params \\
\midrule
\multicolumn{3}{l}{\textit{Query Characteristics}} \\
Query length & 128$\pm$75 & mean$\pm$std (tokens) \\
\quad Range & 13--563 & tokens \\
\quad Median & 105 & tokens \\
\bottomrule
\end{tabular}
\end{table}

\textbf{Tool Diversity.} The 993 unique tools span diverse domains including weather/forecasting, search/retrieval, mathematical computation, messaging, and data access operations. This diversity ensures evaluation across varied function calling scenarios rather than a narrow tool subset.

\textbf{Parameter Complexity.} While 48.7\% of parameters are strings (suitable for semantic comparison), 5.2\% are complex types (lists and dicts) requiring STED's recursive structure handling. As shown in Table~\ref{tab:complexity-correlation}, these complex parameters are disproportionately represented in inconsistent samples (60$\times$ higher for nested structures), justifying STED's structure-aware approach.

\textbf{Multilingual Coverage.} With 56.4\% of queries containing non-ASCII characters spanning Chinese, Japanese, Korean, and European languages, the dataset tests model consistency across linguistic contexts---an important consideration for global deployment.

\subsection{Implementation and Reproducibility Details}
\label{subsec:reproducibility}

\textbf{Embedding Model.} We use Amazon Titan Text Embeddings v2 (\texttt{amazon.titan-embed-text-v2:0}) with 1024-dimensional vectors for semantic similarity computation. For local experiments and validation, we use \texttt{all-MiniLM-L6-v2} (384 dimensions) from Sentence-Transformers, which provides comparable semantic matching quality at lower computational cost. The choice of embedding model has minimal impact on STED's discrimination ability since we primarily compare short field names and values where most modern embedding models achieve similar performance.

\textbf{Hyperparameters.} Default STED configuration: key weight $w = 0.5$ (equal weighting of key similarity and children similarity), stability steepness $\alpha = 20$ (Section~\ref{sec:ablation-steepness}), text chunking threshold 300 characters with 50-character overlap.

\textbf{Random Seeds.} All experiments use seed 42 for reproducibility. For LLM generation, we use each provider's default sampling behavior at specified temperatures.

\textbf{Computational Environment.} Experiments were conducted on AWS using Python 3.11. LLM API calls were made to respective cloud providers (Anthropic, AWS Bedrock, OpenAI, Google, etc.). STED computation runs in approximately 50ms per comparison for typical JSON structures (depth $\leq$ 5, $\leq$ 50 fields) on a single CPU core.

\section{Statistical Analysis Details}
\label{sec:statistical-analysis}

\subsection{Methodology}

We performed comprehensive statistical validation using: \textbf{Paired t-tests}: Each of the 80 base samples was evaluated by all metrics on the same variations, enabling paired comparisons; \textbf{Bonferroni correction}: Applied for multiple comparisons ($\alpha = 0.05/30 = 0.0017$); \textbf{Effect size calculation}: Cohen's d to quantify practical significance; \textbf{Confidence intervals}: 95\% CIs computed using bootstrap with 1000 iterations.

\subsection{Semantic Variations Results}

For semantic variations (N=800): STED (mean 0.9539$\pm$0.0393, 95\% CI [0.9511, 0.9567]), BERTScore (0.9582$\pm$0.0249, [0.9564, 0.9600]), TED (1.0$\pm$0.0, constant), DeepDiff (0.7991$\pm$0.1647, [0.7873, 0.8109]).

STED vs BERTScore: mean p = 0.600 (not significant), demonstrating statistically equivalent performance on semantic variations. STED vs TED and STED vs DeepDiff: p < 0.001*** (significant after Bonferroni correction).

\subsection{Expression Variations Results}

For expression variations (N=800): STED (0.9812$\pm$0.0174, [0.9799, 0.9824]), BERTScore (0.9595$\pm$0.0267, [0.9576, 0.9614]), TED (1.0$\pm$0.0), DeepDiff (0.8051$\pm$0.1644, [0.7934, 0.8169]).

STED significantly outperforms all baselines (p < 0.001), demonstrating superior format-agnostic understanding. Lower standard deviation (0.0174) indicates more consistent performance.

\subsection{Effect Size Analysis}

Cohen's d effect sizes for STED comparisons:
\begin{itemize}
\item STED vs TED: d = -1.23 (semantic), -1.45 (expression) --- Large
\item STED vs BERTScore: d = -0.12 (semantic, negligible), 0.82 (expression) --- Large
\item STED vs DeepDiff: d = 1.26 (semantic), 1.38 (expression) --- Large
\end{itemize}

\section{Detailed Consistency Results}
\label{sec:results}

\subsection{ShareGPT Structured Output Results}
\label{sec:sharegpt-results}

We evaluate 18 models with complete results on 75 valid ShareGPT samples across 11 temperatures (0.0--1.0), yielding 8,250 outputs per model (75 samples $\times$ 10 runs $\times$ 11 temperatures). Table~\ref{tab:consistency-summary} shows aggregated stability metrics, and Table~\ref{tab:overall-consistency-temp} provides stability scores ($S_\alpha$) by temperature.

\begin{table}[h]
\centering
\caption{ShareGPT: Aggregated Stability Score (Mean$\pm$SD across temperatures, 75 samples, 18 models)}
\label{tab:consistency-summary}
\scriptsize
\begin{tabular}{@{}lcccc@{}}
\toprule
\textbf{Model} & \textbf{Overall} & \textbf{Semantic} & \textbf{Structural} & \textbf{Valid.} \\
\midrule
Claude-Sonnet-4.5 & \textbf{.579$\pm$.310} & \textbf{.538} & .782 & 89.5\% \\
Claude-Opus-4.5 & .572$\pm$.303 & .527 & \textbf{.789} & 96.3\% \\
Claude-Opus-4 & .548$\pm$.333 & .509 & .736 & 97.7\% \\
Claude-Haiku-4.5 & .527$\pm$.306 & .465 & .696 & \textbf{100.0\%} \\
Claude-Sonnet-4 & .514$\pm$.289 & .476 & .703 & 89.5\% \\
Claude-3.7-Sonnet & .499$\pm$.292 & .429 & .781 & \textbf{100.0\%} \\
Claude-3.5-Haiku & .496$\pm$.344 & .429 & .734 & 99.1\% \\
Claude-3.5-Sonnet & .486$\pm$.322 & .422 & .765 & \textbf{100.0\%} \\
Gemini-2.5-Flash-Lite & .452$\pm$.347 & .407 & .651 & \textbf{100.0\%} \\
GPT-4.1-Mini & .431$\pm$.321 & .377 & .611 & \textbf{100.0\%} \\
Qwen3-235B-A22B & .421$\pm$.300 & .359 & .634 & \textbf{100.0\%} \\
Qwen3-32B & .362$\pm$.298 & .289 & .639 & 96.8\% \\
Llama-3.3-70B & .336$\pm$.313 & .275 & .649 & 98.3\% \\
Grok-4.1-Fast & .333$\pm$.304 & .267 & .569 & 97.1\% \\
Nova-2-Lite & .325$\pm$.335 & .273 & .519 & \textbf{100.0\%} \\
GPT-OSS-120B & .304$\pm$.281 & .245 & .493 & \textbf{100.0\%} \\
Mimo-V2-Flash & .255$\pm$.285 & .204 & .440 & 74.4\% \\
Minimax-M2 & .226$\pm$.254 & .170 & .420 & 98.2\% \\
\bottomrule
\end{tabular}

\vspace{0.3em}
\raggedright\small Sorted by overall stability. Bold: best in column.
\end{table}

\begin{table*}[t]
\centering
\caption{ShareGPT Structured Output: Stability Score ($S_\alpha$, hybrid $w{=}0.5$) by Temperature (75 samples, 18 models)}
\label{tab:overall-consistency-temp}
\scriptsize
\begin{tabular}{@{}lccccccccccc@{}}
\toprule
\textbf{Model} & \textbf{T=0.0} & \textbf{T=0.1} & \textbf{T=0.2} & \textbf{T=0.3} & \textbf{T=0.4} & \textbf{T=0.5} & \textbf{T=0.6} & \textbf{T=0.7} & \textbf{T=0.8} & \textbf{T=0.9} & \textbf{T=1.0} \\
\midrule
Claude-Sonnet-4.5 & .905 & .664 & .627 & .591 & .574 & .545 & .539 & .486 & .482 & .489 & .471 \\
Claude-Opus-4.5 & .774 & \textbf{.690} & \textbf{.616} & \textbf{.611} & \textbf{.589} & \textbf{.556} & \textbf{.518} & \textbf{.511} & \textbf{.457} & \textbf{.474} & \textbf{.500} \\
Claude-Opus-4 & .934 & .672 & .599 & .564 & .488 & .519 & .488 & .475 & .442 & .440 & .406 \\
Claude-Haiku-4.5 & .965 & .633 & .555 & .553 & .505 & .478 & .444 & .425 & .432 & .420 & .388 \\
Claude-Sonnet-4 & .833 & .582 & .541 & .517 & .505 & .478 & .472 & .430 & .443 & .427 & .428 \\
Claude-3.7-Sonnet & .581 & .573 & .586 & .568 & .529 & .467 & .470 & .429 & .431 & .424 & .428 \\
Claude-3.5-Haiku & .976 & .628 & .551 & .518 & .484 & .409 & .423 & .420 & .408 & .336 & .304 \\
Claude-3.5-Sonnet & .565 & .539 & .530 & .513 & .508 & .494 & .470 & .463 & .445 & .408 & .404 \\
Gemini-2.5-Flash-Lite & \textbf{.987} & .552 & .506 & .437 & .423 & .371 & .368 & .350 & .319 & .347 & .316 \\
GPT-4.1-Mini & .615 & .524 & .498 & .455 & .460 & .389 & .399 & .366 & .372 & .345 & .315 \\
Qwen3-235B-A22B & .508 & .472 & .476 & .463 & .424 & .416 & .374 & .373 & .389 & .384 & .355 \\
Qwen3-32B & .582 & .504 & .443 & .385 & .359 & .342 & .324 & .291 & .280 & .245 & .226 \\
Llama-3.3-70B & .493 & .457 & .390 & .350 & .348 & .299 & .307 & .282 & .248 & .279 & .248 \\
Grok-4.1-Fast & .363 & .332 & .372 & .332 & .342 & .323 & .318 & .315 & .329 & .327 & .305 \\
Nova-2-Lite & .438 & .374 & .389 & .354 & .315 & .310 & .311 & .292 & .277 & .254 & .262 \\
GPT-OSS-120B & .366 & .342 & .295 & .326 & .319 & .310 & .301 & .299 & .286 & .250 & .251 \\
Mimo-V2-Flash & .308 & .308 & .329 & .293 & .269 & .242 & .224 & .239 & .218 & .184 & .188 \\
Minimax-M2 & .344 & .314 & .285 & .288 & .268 & .222 & .209 & .164 & .156 & .124 & .110 \\
\bottomrule
\end{tabular}

\vspace{0.5em}
\raggedright\small Sorted by overall stability. Bold: best in column.
\end{table*}

\begin{table}[h]
\centering
\caption{ShareGPT Structured Output: Validity Rate by Temperature (\%)}
\label{tab:sharegpt-validity-temp}
\resizebox{\columnwidth}{!}{%
\begin{tabular}{@{}lcccccc@{}}
\toprule
\textbf{Model} & \textbf{T=0.0} & \textbf{T=0.2} & \textbf{T=0.4} & \textbf{T=0.6} & \textbf{T=0.8} & \textbf{T=1.0} \\
\midrule
Nova-2-Lite & \textbf{100.0} & \textbf{100.0} & \textbf{100.0} & \textbf{100.0} & \textbf{100.0} & \textbf{100.0} \\
Gemini-2.5-Flash-Lite & \textbf{100.0} & \textbf{100.0} & \textbf{100.0} & \textbf{100.0} & \textbf{100.0} & \textbf{100.0} \\
GPT-4.1-Mini & \textbf{100.0} & \textbf{100.0} & \textbf{100.0} & \textbf{100.0} & \textbf{100.0} & \textbf{100.0} \\
GPT-OSS-120B & \textbf{100.0} & \textbf{100.0} & \textbf{100.0} & \textbf{100.0} & \textbf{100.0} & \textbf{100.0} \\
Claude-Haiku-4.5 & \textbf{100.0} & \textbf{100.0} & \textbf{100.0} & \textbf{100.0} & \textbf{100.0} & \textbf{100.0} \\
Claude-3.7-Sonnet & \textbf{100.0} & \textbf{100.0} & \textbf{100.0} & \textbf{100.0} & \textbf{100.0} & 99.9 \\
Claude-3.5-Sonnet & \textbf{100.0} & \textbf{100.0} & \textbf{100.0} & \textbf{100.0} & \textbf{100.0} & 99.9 \\
Qwen3-235B-A22B & \textbf{100.0} & \textbf{100.0} & 99.9 & \textbf{100.0} & \textbf{100.0} & \textbf{100.0} \\
Claude-3.5-Haiku & 99.6 & 99.0 & 99.6 & 99.2 & 98.4 & 99.9 \\
Llama-3.3-70B & 99.6 & 97.8 & 99.4 & 98.0 & 98.0 & 96.9 \\
Grok-4.1-Fast & 98.1 & 97.9 & 97.9 & 98.5 & 94.6 & 94.4 \\
Minimax-M2 & 96.6 & 95.0 & 99.2 & 98.5 & 98.5 & 99.4 \\
Claude-Opus-4 & 96.2 & 98.4 & 98.5 & 98.1 & 97.4 & 98.5 \\
Qwen3-32B & 96.2 & 96.6 & 96.8 & 97.2 & 97.2 & 97.1 \\
Claude-Opus-4.5 & 96.4 & 96.8 & 96.0 & 95.9 & 95.2 & 96.6 \\
Claude-Sonnet-4.5 & 91.0 & 87.2 & 89.4 & 89.2 & 91.2 & 92.1 \\
Mimo-V2-Flash & 74.9 & 74.1 & 74.5 & 74.1 & 74.1 & 74.2 \\
Claude-Sonnet-4 & 72.9 & 69.8 & 70.6 & 71.8 & 72.5 & 72.2 \\
\bottomrule
\end{tabular}%
}

\vspace{0.3em}
\raggedright\scriptsize Sorted by avg validity. Bold: 100\%. 18 models evaluated on 75 samples.
\end{table}

\textbf{Structural vs Semantic Consistency.} For structured output tasks, structural consistency (schema adherence) is critical for downstream processing, while semantic consistency (content similarity) reflects value stability. Tables~\ref{tab:sharegpt-structural-temp} and~\ref{tab:sharegpt-semantic-temp} decompose STED into these components.

\begin{table*}[t]
\centering
\caption{ShareGPT Structured Output: Structural Stability ($S_\alpha$, $w{=}1.0$) by Temperature (75 samples, 18 models)}
\label{tab:sharegpt-structural-temp}
\scriptsize
\begin{tabular}{@{}lccccccccccc@{}}
\toprule
\textbf{Model} & \textbf{T=0.0} & \textbf{T=0.1} & \textbf{T=0.2} & \textbf{T=0.3} & \textbf{T=0.4} & \textbf{T=0.5} & \textbf{T=0.6} & \textbf{T=0.7} & \textbf{T=0.8} & \textbf{T=0.9} & \textbf{T=1.0} \\
\midrule
Gemini-2.5-Flash-Lite & \textbf{.988} & .716 & .685 & .640 & .636 & .612 & .611 & .583 & .562 & .593 & .542 \\
Claude-3.5-Haiku & .988 & .814 & .806 & .778 & .749 & .704 & .718 & .737 & .705 & .587 & .574 \\
Claude-Haiku-4.5 & .982 & .748 & .706 & .712 & .670 & .665 & .646 & .633 & .636 & .634 & .633 \\
Claude-Opus-4 & .967 & .796 & .726 & .752 & .681 & .730 & .708 & .698 & .678 & .699 & .666 \\
Claude-Sonnet-4.5 & .946 & .844 & .816 & .764 & .792 & .778 & .753 & .740 & .726 & .735 & .706 \\
Claude-Opus-4.5 & .891 & \textbf{.884} & \textbf{.832} & \textbf{.846} & \textbf{.823} & \textbf{.830} & \textbf{.781} & \textbf{.807} & \textbf{.758} & \textbf{.770} & \textbf{.784} \\
Claude-Sonnet-4 & .821 & .809 & .811 & .718 & .739 & .720 & .749 & .710 & .720 & .706 & .693 \\
Claude-3.5-Sonnet & .806 & .787 & .795 & .787 & .793 & .768 & .775 & .753 & .729 & .717 & .712 \\
Claude-3.7-Sonnet & .802 & .795 & .843 & .842 & .808 & .755 & .795 & .735 & .766 & .730 & .731 \\
Qwen3-32B & .789 & .743 & .731 & .695 & .665 & .657 & .614 & .566 & .572 & .533 & .514 \\
Llama-3.3-70B & .769 & .745 & .711 & .676 & .673 & .656 & .664 & .614 & .593 & .605 & .576 \\
GPT-4.1-Mini & .753 & .696 & .667 & .640 & .636 & .578 & .574 & .571 & .570 & .549 & .494 \\
Qwen3-235B-A22B & .690 & .688 & .674 & .652 & .618 & .651 & .598 & .593 & .619 & .617 & .580 \\
Grok-4.1-Fast & .602 & .588 & .608 & .552 & .589 & .603 & .578 & .567 & .587 & .607 & .621 \\
Minimax-M2 & .598 & .520 & .522 & .540 & .514 & .405 & .431 & .385 & .353 & .266 & .243 \\
Nova-2-Lite & .586 & .523 & .579 & .540 & .566 & .539 & .520 & .479 & .461 & .444 & .479 \\
GPT-OSS-120B & .564 & .540 & .466 & .513 & .500 & .505 & .481 & .502 & .465 & .457 & .431 \\
Mimo-V2-Flash & .491 & .491 & .517 & .455 & .474 & .452 & .449 & .428 & .443 & .421 & .398 \\
\bottomrule
\end{tabular}

\vspace{0.5em}
\raggedright\small Sorted by T=0.0 structural stability. Bold: best in column.
\end{table*}

\begin{table*}[t]
\centering
\caption{ShareGPT Structured Output: Semantic Stability ($S_\alpha$, $w{=}0.0$) by Temperature (75 samples, 18 models)}
\label{tab:sharegpt-semantic-temp}
\scriptsize
\begin{tabular}{@{}lccccccccccc@{}}
\toprule
\textbf{Model} & \textbf{T=0.0} & \textbf{T=0.1} & \textbf{T=0.2} & \textbf{T=0.3} & \textbf{T=0.4} & \textbf{T=0.5} & \textbf{T=0.6} & \textbf{T=0.7} & \textbf{T=0.8} & \textbf{T=0.9} & \textbf{T=1.0} \\
\midrule
Gemini-2.5-Flash-Lite & \textbf{.989} & .533 & .465 & .395 & .369 & .336 & .316 & .290 & .257 & .285 & .250 \\
Claude-3.5-Haiku & .981 & .582 & .481 & .452 & .415 & .345 & .342 & .331 & .329 & .266 & .242 \\
Claude-Haiku-4.5 & .970 & .582 & .492 & .468 & .425 & .415 & .391 & .364 & .360 & .339 & .310 \\
Claude-Opus-4 & .943 & \textbf{.673} & .586 & .528 & .446 & .458 & .430 & .418 & .385 & .372 & .361 \\
Claude-Sonnet-4.5 & .902 & .653 & .586 & .553 & .529 & .511 & .496 & .433 & .419 & .431 & .404 \\
Claude-Opus-4.5 & .767 & .706 & \textbf{.620} & \textbf{.594} & \textbf{.567} & \textbf{.521} & \textbf{.477} & \textbf{.490} & \textbf{.448} & \textbf{.452} & \textbf{.460} \\
Claude-Sonnet-4 & .626 & .596 & .592 & .484 & .487 & .464 & .476 & .411 & .435 & .440 & .393 \\
GPT-4.1-Mini & .590 & .474 & .465 & .400 & .407 & .337 & .325 & .311 & .305 & .282 & .258 \\
Claude-3.7-Sonnet & .540 & .524 & .535 & .490 & .462 & .402 & .395 & .362 & .349 & .335 & .331 \\
Claude-3.5-Sonnet & .517 & .503 & .498 & .448 & .452 & .418 & .404 & .397 & .370 & .325 & .319 \\
Qwen3-32B & .514 & .434 & .372 & .304 & .300 & .282 & .238 & .219 & .208 & .166 & .156 \\
Llama-3.3-70B & .446 & .414 & .337 & .303 & .287 & .236 & .257 & .232 & .194 & .202 & .169 \\
Qwen3-235B-A22B & .445 & .398 & .415 & .388 & .355 & .364 & .317 & .303 & .348 & .321 & .304 \\
Nova-2-Lite & .399 & .341 & .339 & .294 & .264 & .237 & .251 & .238 & .228 & .215 & .201 \\
GPT-OSS-120B & .298 & .299 & .232 & .279 & .264 & .247 & .228 & .243 & .224 & .199 & .190 \\
Grok-4.1-Fast & .291 & .264 & .308 & .278 & .286 & .257 & .254 & .269 & .303 & .300 & .254 \\
Minimax-M2 & .274 & .270 & .265 & .223 & .194 & .165 & .157 & .118 & .108 & .081 & .076 \\
Mimo-V2-Flash & .272 & .260 & .279 & .244 & .220 & .201 & .169 & .196 & .156 & .131 & .140 \\
\bottomrule
\end{tabular}

\vspace{0.5em}
\raggedright\small Sorted by T=0.0 semantic stability. Bold: best in column.
\end{table*}

\textbf{Key Insight:} Structural consistency is substantially higher than semantic consistency (mean 0.68 vs 0.41 at T=0.0), indicating models maintain schema structure more reliably than content values.

\textbf{Key Observations:} Claude-Opus-4.5 leads overall consistency (0.676$\pm$0.307), with Claude-3.7-Sonnet achieving highest structural consistency (0.788). Four models achieve 100\% validity: Claude-3.7-Sonnet, Claude-3.5-Sonnet, Claude-Haiku-4.5, and Qwen3-235B-A22B. Temperature sensitivity varies widely: Claude-3.5-Haiku shows strongest decline (28.2\% drop), while Grok-4.1-Fast exhibits exceptional stability (only 0.9\% decline). Llama-3.3-70B achieves high validity (98.4\%) on structured output but only 32.7\% on tool calling, demonstrating critical task-specific differences.

\begin{table}[h]
\centering
\caption{Prompt Factors vs Consistency (ShareGPT, All Temperatures, N=77)}
\label{tab:sharegpt-complexity-correlation}
\scriptsize
\begin{tabular}{@{}lccc@{}}
\toprule
\textbf{Factor} & \textbf{Consistent ($\geq$0.51)} & \textbf{Inconsistent ($\leq$0.35)} & \textbf{Ratio} \\
\midrule
\multicolumn{4}{l}{\textit{Prompt characteristics}} \\
Question marks (avg) & 0.8 & 2.4 & 2.8$\times$ \\
Prompt length (chars) & 15,298 & 15,877 & 1.0$\times$ \\
Prose-heavy$^a$ & 5.0\% & 5.0\% & 1.0$\times$ \\
\midrule
Sample count & 20 (26\%) & 20 (26\%) & -- \\
\bottomrule
\end{tabular}

\vspace{0.3em}
\raggedright\scriptsize $^a$Avg sentence $>$20 words, no structure markers. Quartile-based thresholds (Q1, Q3) across 18 models, all temperatures. Correlation (constraint score vs stability): $r = 0.03$, $p = 0.80$ (not significant).
\end{table}

\subsection{Toucan Tool Calling Results}
\label{sec:toucan-full-results}

We evaluate 18 final models on 1,006 Toucan samples across 11 temperatures (0.0--1.0), yielding 110,660 outputs per model (1,006 samples $\times$ 10 runs $\times$ 11 temperatures). Table~\ref{tab:toucan-full-temp} presents stability scores ($S_\alpha$) by temperature for all 18 models, extending the 12-model summary in main paper Table~\ref{tab:toucan-stability}. Table~\ref{tab:toucan-validity-temp} shows validity rates.

\begin{table*}[t]
\centering
\caption{Toucan Tool Calling: Stability Score ($S_\alpha$) by Temperature (1,006 samples, 18 models)}
\label{tab:toucan-full-temp}
\scriptsize
\begin{tabular}{@{}lccccccccccc@{}}
\toprule
\textbf{Model} & \textbf{T=0.0} & \textbf{T=0.1} & \textbf{T=0.2} & \textbf{T=0.3} & \textbf{T=0.4} & \textbf{T=0.5} & \textbf{T=0.6} & \textbf{T=0.7} & \textbf{T=0.8} & \textbf{T=0.9} & \textbf{T=1.0} \\
\midrule
Claude-3.5-Haiku & \textbf{.970} & \textbf{.905} & \textbf{.886} & \textbf{.872} & .857 & \textbf{.854} & .841 & .839 & .826 & .827 & \textbf{.826} \\
Claude-Opus-4 & .969 & .941 & .911 & .899 & .883 & .862 & .854 & .830 & .827 & .802 & .802 \\
Claude-Sonnet-4.5 & .959 & .883 & .842 & .814 & .809 & .789 & .778 & .766 & .738 & .736 & .719 \\
Claude-Sonnet-4 & .959 & .879 & .839 & .810 & .804 & .785 & .773 & .760 & .732 & .732 & .714 \\
Claude-Opus-4.5 & .948 & .902 & .863 & .845 & .828 & .807 & .791 & .763 & .765 & .750 & .744 \\
Claude-3.7-Sonnet & .891 & .879 & .872 & .869 & \textbf{.858} & \textbf{.854} & \textbf{.852} & \textbf{.840} & \textbf{.843} & \textbf{.834} & .824 \\
Claude-3.5-Sonnet & .888 & .880 & .873 & .853 & .851 & .830 & .839 & .827 & .817 & .804 & .808 \\
Nova-2-Lite & .869 & .857 & .852 & .849 & .851 & .851 & .832 & .828 & .830 & .821 & .814 \\
Claude-Haiku-4.5 & .920 & .823 & .764 & .729 & .701 & .679 & .648 & .637 & .627 & .613 & .579 \\
Qwen3-235B-A22B & .746 & .747 & .743 & .723 & .725 & .712 & .688 & .666 & .646 & .628 & .608 \\
Mimo-V2-Flash & .761 & .745 & .730 & .718 & .689 & .657 & .640 & .618 & .595 & .572 & .534 \\
Qwen3-32B & .703 & .695 & .675 & .661 & .630 & .624 & .598 & .578 & .545 & .526 & .513 \\
Grok-4.1-Fast & .588 & .562 & .567 & .563 & .575 & .551 & .562 & .558 & .541 & .531 & .548 \\
GPT-4.1-Mini & .542 & .531 & .536 & .521 & .553 & .593 & .578 & .565 & .566 & .547 & .542 \\
GPT-OSS-120B & .587 & .570 & .573 & .555 & .567 & .542 & .523 & .529 & .503 & .498 & .505 \\
Minimax-M2 & .504 & .501 & .477 & .464 & .437 & .443 & .404 & .397 & .386 & .350 & .330 \\
Gemini-2.5-Flash-Lite & .540 & .473 & .448 & .429 & .343 & .382 & .380 & .362 & .329 & .327 & .293 \\
Llama-3.3-70B & .306 & .305 & .300 & .295 & .302 & .295 & .290 & .287 & .290 & .279 & .279 \\
\bottomrule
\end{tabular}

\vspace{0.5em}
\raggedright\small Higher = more stable. Sorted by T=0.0 stability. Full results for 18 models; main Table~\ref{tab:toucan-stability} shows 12 representative models.
\end{table*}

\begin{table}[h]
\centering
\caption{Toucan Tool Calling: Validity Rate by Temperature (\%)}
\label{tab:toucan-validity-temp}
\resizebox{\columnwidth}{!}{%
\begin{tabular}{@{}lcccccc@{}}
\toprule
\textbf{Model} & \textbf{T=0.0} & \textbf{T=0.2} & \textbf{T=0.4} & \textbf{T=0.6} & \textbf{T=0.8} & \textbf{T=1.0} \\
\midrule
Qwen3-235B-A22B & \textbf{98.1} & \textbf{98.1} & \textbf{98.1} & \textbf{98.1} & \textbf{98.1} & \textbf{98.1} \\
Mimo-V2-Flash & 98.3 & 98.0 & 98.0 & 97.8 & 97.3 & 96.8 \\
Claude-3.5-Haiku & 97.7 & 97.7 & 97.7 & 97.7 & 97.7 & 97.7 \\
Claude-3.5-Sonnet & 97.5 & 97.4 & 97.4 & 97.4 & 97.5 & 97.4 \\
Claude-Opus-4.5 & 97.7 & 97.7 & 97.6 & 97.5 & 95.7 & 95.1 \\
Claude-Sonnet-4 & 97.2 & 97.1 & 97.0 & 97.0 & 97.0 & 97.0 \\
Claude-Opus-4 & 97.0 & 96.9 & 97.1 & 97.1 & 97.2 & 97.2 \\
Claude-Sonnet-4.5 & 96.9 & 96.8 & 96.9 & 96.8 & 96.8 & 96.8 \\
Claude-3.7-Sonnet & 96.7 & 96.7 & 96.6 & 96.7 & 96.7 & 96.8 \\
Qwen3-32B & 96.6 & 96.8 & 96.6 & 96.6 & 96.5 & 96.4 \\
Nova-2-Lite & 96.6 & 96.5 & 96.6 & 96.5 & 96.6 & 96.4 \\
Claude-Haiku-4.5 & 95.3 & 95.4 & 95.2 & 95.2 & 95.3 & 95.0 \\
Grok-4.1-Fast & 95.0 & 94.7 & 94.7 & 94.8 & 95.0 & 95.0 \\
Minimax-M2 & 90.9 & 90.6 & 90.4 & 90.7 & 90.6 & 90.3 \\
GPT-OSS-120B & 82.4 & 82.3 & 82.6 & 82.2 & 82.2 & 81.8 \\
GPT-4.1-Mini & 75.8 & 76.0 & 76.0 & 76.0 & 75.9 & 75.9 \\
Gemini-2.5-Flash-Lite & 54.2 & 54.3 & 45.3 & 53.3 & 53.7 & 52.4 \\
Llama-3.3-70B & 32.5 & 32.0 & 32.7 & 32.6 & 32.9 & 32.5 \\
\bottomrule
\end{tabular}%
}

\vspace{0.3em}
\raggedright\scriptsize Sorted by avg validity. Bold: highest. 18 models evaluated on 1,006 samples.
\end{table}

\begin{table}[h]
\centering
\caption{Toucan Tool Calling: Inconsistency Pattern Analysis at T=0.5 (1,006 samples)}
\label{tab:inconsistency-patterns}
\resizebox{\columnwidth}{!}{%
\begin{tabular}{@{}lcccccc@{}}
\toprule
\textbf{Model} & \textbf{Consist.} & \textbf{Tool Sel.} & \textbf{Param.} & \textbf{Count} & \textbf{Malf.} & \textbf{Uniq} \\
\midrule
Llama-3.3-70B & 88.0\% & 1.8\% & 10.0\% & 0.0\% & 0.2\% & 1.38 \\
Claude-Opus-4 & 74.8\% & 7.6\% & 16.3\% & 5.3\% & 1.5\% & 2.05 \\
Nova-2-Lite & 73.7\% & 6.8\% & 18.5\% & 5.6\% & 1.5\% & 2.05 \\
Claude-3.5-Haiku & 72.9\% & 7.9\% & 17.6\% & 5.1\% & 1.3\% & 2.16 \\
Claude-3.7-Sonnet & 71.8\% & 7.2\% & 20.4\% & 0.0\% & 1.3\% & 2.10 \\
Claude-3.5-Sonnet & 71.5\% & 10.8\% & 16.0\% & 6.4\% & 1.3\% & 2.18 \\
Gemini-2.5-Flash-Lite & 71.6\% & 16.4\% & 10.0\% & 13.4\% & 1.5\% & 1.72 \\
Claude-Sonnet-4 & 70.0\% & 11.0\% & 17.5\% & 9.5\% & 1.9\% & 2.18 \\
Claude-Haiku-4.5 & 60.7\% & 18.2\% & 16.9\% & 20.7\% & 2.1\% & 2.50 \\
Qwen3-235B-A22B & 59.7\% & 17.4\% & 17.5\% & 20.2\% & 1.9\% & 2.72 \\
Claude-Opus-4.5 & 67.2\% & 11.5\% & 17.4\% & 13.4\% & 2.2\% & 2.46 \\
Claude-Sonnet-4.5 & 67.1\% & 11.7\% & 16.8\% & 13.8\% & 2.0\% & 2.41 \\
Mimo-V2-Flash & 55.0\% & 24.2\% & 18.8\% & 20.8\% & 1.9\% & 2.61 \\
Qwen3-32B & 48.0\% & 14.1\% & 36.0\% & 11.0\% & 1.5\% & 2.88 \\
Minimax-M2 & 41.0\% & 37.9\% & 19.1\% & 30.1\% & 2.2\% & 3.33 \\
\bottomrule
\end{tabular}%
}

\vspace{0.3em}
\raggedright\scriptsize Consist.: identical outputs across all runs. Tool Sel.: different tools. Param.: different parameter values. Count: different tool call count. Malf.: parsing errors. Uniq: avg unique outputs.
\end{table}

\subsection{Summary and Recommendations}

\textbf{Temperature Selection:} T=0.1--0.3 provides optimal consistency for both tasks. T=0.4--0.6 offers acceptable trade-off between diversity and consistency. T=0.7--0.9 significantly degrades consistency, especially semantic preservation.

\textbf{Model Selection:} For structured output requiring high consistency, Claude-3.7-Sonnet excels on both tasks. For tool calling specifically, Claude-Opus-4 achieves highest peak consistency (0.948). For validity-critical applications, Qwen3-235B-A22B leads (98.1\% on Toucan, 100\% on ShareGPT).

\textbf{Task-Specific Considerations:} Models exhibit dramatically different performance across tasks---Llama-3.3-70B achieves 99\% validity on structured output but only 33\% on tool calling. Always evaluate on task-specific benchmarks rather than relying on general capabilities.

\begin{table}[h]
\centering
\caption{Parameter Complexity Correlation with Inconsistency (Toucan, T=0.5)}
\label{tab:complexity-correlation}
\scriptsize
\begin{tabular}{@{}lccc@{}}
\toprule
\textbf{Parameter Type} & \textbf{Consistent} & \textbf{Inconsistent} & \textbf{Ratio} \\
\midrule
Has Nested Structure & 0.1\% & 6.4\% & 60$\times$ \\
Has Dict Params & 0.6\% & 7.5\% & 12.5$\times$ \\
Has List Params & 3.5\% & 12.1\% & 3.4$\times$ \\
Has Any Complex Type & 4.1\% & 16.4\% & 4.0$\times$ \\
\midrule
Mean Max Depth & 0.04 & 0.27 & 6.8$\times$ \\
\bottomrule
\end{tabular}

\vspace{0.5em}
\raggedright\small Analysis across samples at T=0.5. Complex parameters are dramatically over-represented in inconsistent samples, justifying STED's structure-aware comparison.
\end{table}

\subsection{Consistency Diagnostics: Failure Pattern Taxonomy}
\label{sec:sted-value-analysis}

A central contribution of our diagnostics framework is characterizing \textit{how} LLM outputs fail to be consistent. While tool calling responses have a fixed structure (tool name + JSON arguments), analysis of inconsistent samples across models reveals four distinct failure categories:

\textbf{1. String Parameter Variations (39.3\%):} The largest category involves semantically equivalent string values. STED's embedding-based comparison recognizes equivalence (e.g., \texttt{"matrix\_A"} $\approx$ \texttt{"matrix\_2s"}, \texttt{"user\_email"} $\approx$ \texttt{"email\_address"}), unlike exact matching which scores 0.

\textbf{2. Tool Count Variations (35.4\%):} Models often generate different numbers of tool calls for the same query. STED handles uneven comparisons through insert/delete costs, providing proportional similarity rather than binary failure.

\textbf{3. Tool Selection Variations (15.8\%):} Different tool subsets are chosen across runs. STED's Hungarian matching provides order-invariant comparison and proportional credit for partial overlap.

\textbf{4. Complex Parameter Variations (5.6\%):} Lists and dictionaries in arguments require recursive structure comparison. While only 5.6\% of inconsistent samples, these are disproportionately represented: samples with nested structures are 60$\times$ more likely to be inconsistent (Table~\ref{tab:complexity-correlation}).

\textbf{Parameter Type Distribution in Inconsistent Samples:} Among 4,631 differing parameters: strings 61.5\%, lists 13.9\%, integers 11.7\%, floats 6.8\%, dicts 3.5\%, booleans 2.5\%. By complexity: primitives 82.5\%, simple lists 10.8\%, nested structures 5.5\%, simple dicts 1.1\%.

The key insight is that while most differing parameters are primitives, \textbf{samples containing complex structures are dramatically more likely to exhibit inconsistency}. This concentration of inconsistency in complex cases justifies STED's structure-aware, semantic approach over simpler exact-matching methods.

\section{Implementation Optimizations}
\label{sec:optimizations}

We describe the key optimizations that reduce STED's time complexity from exponential to polynomial while maintaining accuracy.

\subsection{Complexity Reduction}

\begin{table}[h]
\centering
\caption{Time Complexity Reduction Through Optimizations}
\label{tab:complexity-reduction}
\scriptsize
\begin{tabular}{@{}lcc@{}}
\toprule
\textbf{Component} & \textbf{Naive} & \textbf{Optimized} \\
\midrule
Subtree comparisons & $O(B^{2D})$ & $O(B^2 \times D)$ \\
Per-level matching & $O(B^3)$ & $O(B^3)$ \\
Combined variation & 2 passes & 1 pass \\
\midrule
\textbf{Total} & $O(B^{2D+3})$ & $O(B^3 \times D)$ \\
\bottomrule
\end{tabular}

\vspace{0.3em}
\raggedright\small $B$: max branching factor, $D$: tree depth. For typical JSON ($B<50$, $D<10$), this reduces from $>10^{30}$ to $<10^{8}$ operations.
\end{table}

\subsection{Optimization Techniques}

\textbf{1. LRU Memoization Cache.} Subtree comparisons are memoized using content-based hashes as keys. An LRU eviction policy bounds memory. Cache keys are tuples of subtree hashes; hits return cached results directly. This eliminates the exponential blowup from repeated subtree comparisons, reducing $O(B^{2D})$ to $O(B^2 \times D)$ unique comparisons.

\textbf{2. Single-Pass Combined Matching.} For ``combined'' variation type (structural + content), we build both cost matrices in a single traversal over all child pairs, computing structural and content costs simultaneously. This halves the tree traversal cost compared to separate passes.

\textbf{3. Lazy Hash Computation.} Node hashes for cache keys are computed lazily and cached by Python object ID, avoiding redundant hash computations during recursive traversal.

\textbf{4. Embedding Precomputation.} The dominant cost is embedding computation (85\% of runtime on cache miss). We precompute embeddings for:
\begin{itemize}
\item Raw field names and values from all samples
\item Preprocessed variants (e.g., \texttt{``userName''} $\rightarrow$ \texttt{``user name''})
\end{itemize}
This reduces embedding cache misses from 40\% to $<$1\%.

\textbf{5. Optional: Greedy Matching.} For very large arrays ($B>100$), Hungarian's $O(B^3)$ becomes a bottleneck. An optional greedy approximation achieves $O(B^2)$ by iteratively matching each child to its best unmatched partner. Trade-off: 30$\times$ speedup but correlation drops from 0.96 to 0.70.

\subsection{Benchmark Results}

\begin{table}[h]
\centering
\caption{STED Performance Benchmarks (50 Toucan sample pairs)}
\label{tab:optimization-benchmarks}
\scriptsize
\begin{tabular}{@{}lccc@{}}
\toprule
\textbf{Configuration} & \textbf{Time (s)} & \textbf{Speedup} & \textbf{Correlation} \\
\midrule
Naive (no cache) & 1.035 & 1.0$\times$ & 1.000 \\
Optimized (cold cache) & 0.148 & 7.0$\times$ & 0.979 \\
Optimized (warm cache) & 0.046 & 22.5$\times$ & 0.979 \\
Greedy mode & 0.015 & 69$\times$ & 0.704 \\
\bottomrule
\end{tabular}

\vspace{0.3em}
\raggedright\small Benchmarks on typical JSON structures (depth 4--7, 20--100 fields). Cold cache: first run; warm cache: repeated comparisons with precomputed embeddings.
\end{table}

\textbf{Accuracy Analysis.} The optimized algorithm maintains high fidelity:
\begin{itemize}
\item Mean Absolute Error vs naive: 0.027
\item Max Absolute Error: 0.284 (rare edge cases)
\item Pearson correlation: 0.979
\end{itemize}

The small accuracy difference stems from hash collisions in the memoization cache (different subtrees with identical hashes). \textbf{Important: All benchmark results reported in this paper use the exact (naive) implementation without caching optimizations}, ensuring collision-free computation. The optimizations are provided for production deployment where the speed-accuracy tradeoff is acceptable. For evaluation/benchmarking applications requiring exact reproducibility, we recommend disabling the subtree cache (\texttt{subtree\_cache\_size=0}).

\subsection{Configuration Options}

Key parameters: \texttt{subtree\_cache\_size} (default 10K, bounds LRU cache), \texttt{use\_greedy\_matching} (enables $O(B^2)$ approximation), \texttt{use\_fp16\_embeddings} (50\% memory reduction), \texttt{early\_pruning\_threshold} (skip threshold, default 0.8).

\textbf{Recommendations:} (1) \textit{Production}: default settings; (2) \textit{Large datasets}: cache 50K, fp16 enabled; (3) \textit{Speed-critical}: greedy matching (30\% accuracy loss); (4) \textit{Memory-constrained}: cache 5K, fp16 enabled.

\section{Stability Score Analysis}
\label{sec:stability-analysis}

We analyze how the stability score $S_{\alpha}$ varies across models and temperatures. The $\alpha$ ablation study is presented in Section~\ref{sec:ablation-steepness} of the main paper.

Figure~\ref{fig:appendix-stability} provides a comprehensive view of stability behavior across 18 LLMs on ShareGPT structured output tasks.

\begin{figure*}[t]
\centering
\includegraphics[width=0.95\textwidth]{figures/appendix_stability_distribution.png}
\caption{Stability score analysis across LLMs on ShareGPT structured output tasks. (a) Box plots showing the distribution of $S_{\alpha}$ for all 18 evaluated models, ordered by median stability. The gradient coloring (red to green) reflects ranking. (b) Temperature degradation curves for five representative models: Claude-Opus-4.5 (most capable proprietary), Qwen3-235B-A22B (most capable open-source), GPT-4.1-Mini (cost-effective), Grok-4.1-Fast (fast inference), and Nova-2-Lite (AWS). Higher temperature generally reduces stability, though degradation patterns vary by model architecture.}
\label{fig:appendix-stability}
\end{figure*}

Key observations from Figure~\ref{fig:appendix-stability}:

\begin{itemize}
\item \textbf{Claude models dominate stability rankings}: Claude-Opus-4.5 achieves the highest median stability, followed by other Claude variants (3.5-Sonnet, Haiku-4.5, 3.5-Haiku).
\item \textbf{Temperature sensitivity varies by model}: Claude-Opus-4.5 shows gradual degradation from $S_{\alpha} \approx 0.77$ (T=0.0) to $\approx 0.50$ (T=1.0), while Nova-2-Lite remains relatively flat across temperatures.
\item \textbf{Open-source models show moderate stability}: Qwen3-235B-A22B and GPT-4.1-Mini achieve competitive stability scores, demonstrating that open-source models can match proprietary ones on this metric.
\end{itemize}

\section{Prompt Characteristic Analysis}
\label{sec:structural-ambiguity-analysis}

We investigate what characteristics of input prompts are associated with output inconsistency in structured generation tasks.

\subsection{Methodology}

We analyze 77 ShareGPT samples across 18 final models and all temperatures (0.0--1.0), computing mean stability score per sample and using quartile-based classification: \textit{consistent} ($S_\alpha \geq 0.51$, top quartile, N=20) and \textit{inconsistent} ($S_\alpha \leq 0.35$, bottom quartile, N=20).

\textbf{Feature Extraction.} We extract structural and textual features from prompts:

\begin{itemize}
\item \textbf{Question marks}: Count of ``?'' characters in prompt
\item \textbf{Prompt length}: Character count
\item \textbf{Prose-heavy}: Average sentence length $>$20 words AND no structure markers
\item \textbf{JSON examples}: Explicit structure in prompt (regex: \texttt{\textbackslash\{.*".*":})
\end{itemize}

\subsection{Results}

\textbf{Statistical Analysis.} No strong correlations were found between prompt characteristics and consistency:
\begin{itemize}
\item Constraint score correlation: $r = 0.03$, $p = 0.80$ (not significant)
\item Prompt length correlation: $r = -0.01$, $p = 0.96$ (not significant)
\end{itemize}

\textbf{Feature Comparison.} Table~\ref{tab:structural-ambiguity-detailed} shows the full breakdown:

\begin{table}[h]
\centering
\caption{Prompt Characteristic Analysis: Full Results (All Temperatures)}
\label{tab:structural-ambiguity-detailed}
\scriptsize
\begin{tabular}{@{}lccc@{}}
\toprule
\textbf{Feature} & \textbf{Consistent} & \textbf{Inconsistent} & \textbf{Ratio} \\
\midrule
Question marks (avg) & 0.8 & 2.4 & 2.8$\times$ \\
Prompt length (chars) & 15,298 & 15,877 & 1.0$\times$ \\
Prose-heavy & 5.0\% & 5.0\% & 1.0$\times$ \\
\bottomrule
\end{tabular}
\end{table}

\subsection{Interpretation}

The analysis reveals that \textbf{prompt characteristics have limited predictive power} for consistency in this dataset. The most notable difference---question mark frequency (2.8$\times$)---suggests prompts with multiple clarification requests may introduce some ambiguity, but this does not reach statistical significance.

\textbf{Key Finding.} Consistency appears to be primarily driven by model capabilities and inherent task difficulty rather than easily measurable prompt characteristics. This suggests that improving consistency requires model-level interventions (e.g., model selection, temperature tuning) rather than prompt-level modifications alone.

\textbf{Practical Implications.}
\begin{enumerate}
\item Model selection and temperature tuning have larger impact than prompt engineering for consistency
\item Use lower temperatures (T=0.0--0.3) for consistency-critical applications
\item Consider parameter complexity (see Toucan analysis) as a stronger predictor of inconsistency
\end{enumerate}

\end{document}
